\haddockmoduleheading{Data.List}
\label{module:Data.List}
\haddockbeginheader
{\haddockverb\begin{verbatim}
module Data.List (
    (++), head, last, tail, init, null, length, map, reverse, intersperse,
    intercalate, transpose, subsequences, permutations, foldl, foldl', foldl1,
    foldl1', foldr, foldr1, concat, concatMap, and, or, any, all, sum, product,
    maximum, minimum, scanl, scanl1, scanr, scanr1, mapAccumL, mapAccumR,
    iterate, repeat, replicate, cycle, unfoldr, take, drop, splitAt, takeWhile,
    dropWhile, span, break, stripPrefix, group, inits, tails, isPrefixOf,
    isSuffixOf, isInfixOf, elem, notElem, lookup, find, filter, partition, (!!),
    elemIndex, elemIndices, findIndex, findIndices, zip, zip3, zip4, zip5, zip6,
    zip7, zipWith, zipWith3, zipWith4, zipWith5, zipWith6, zipWith7, unzip,
    unzip3, unzip4, unzip5, unzip6, unzip7, lines, words, unlines, unwords, nub,
    delete, (\\), union, intersect, sort, insert, nubBy, deleteBy,
    deleteFirstsBy, unionBy, intersectBy, groupBy, sortBy, insertBy, maximumBy,
    minimumBy, genericLength, genericTake, genericDrop, genericSplitAt,
    genericIndex, genericReplicate
  ) where\end{verbatim}}
\haddockendheader

\section{Basic functions}
\begin{haddockdesc}
\item[\begin{tabular}{@{}l}
(++) :: {\char 91}a{\char 93} -> {\char 91}a{\char 93} -> {\char 91}a{\char 93}
\end{tabular}]
{\haddockbegindoc
\haddockid{(++)} appends two lists, i.e.,\par
\begin{quote}
{\haddockverb\begin{verbatim}
[x1, ..., xm] ++ [y1, ..., yn] == [x1, ..., xm, y1, ..., yn]
[x1, ..., xm] ++ [y1, ...] == [x1, ..., xm, y1, ...]\end{verbatim}}
\end{quote}
If the first list is not finite, the result is the first list.\par
\subsubsection*{\textbf{Performance considerations}}
This function takes linear time in the number of elements of the
 \textbf{first} list. Thus it is better to associate repeated
 applications of \haddockid{(++)} to the right (which is the default behaviour):
 \haddocktt{xs ++ (ys ++ zs)} or simply \haddocktt{xs ++ ys ++ zs}, but not \haddocktt{(xs ++ ys) ++ zs}.
 For the same reason \haddockid{concat} \haddocktt{=} \haddockid{foldr} \haddockid{(++)} \haddocktt{{\char 91}{\char 93}}
 has linear performance, while \haddockid{foldl} \haddockid{(++)} \haddocktt{{\char 91}{\char 93}} is prone
 to quadratic slowdown\par
\subsubsection*{\textbf{Examples}}
\begin{quote}
{\haddockverb\begin{verbatim}
>>> [1, 2, 3] ++ [4, 5, 6]
[1,2,3,4,5,6]

\end{verbatim}}
\end{quote}
\begin{quote}
{\haddockverb\begin{verbatim}
>>> [] ++ [1, 2, 3]
[1,2,3]

\end{verbatim}}
\end{quote}
\begin{quote}
{\haddockverb\begin{verbatim}
>>> [3, 2, 1] ++ []
[3,2,1]

\end{verbatim}}
\end{quote}}
\end{haddockdesc}
\begin{haddockdesc}
\item[\begin{tabular}{@{}l}
head :: HasCallStack => {\char 91}a{\char 93} -> a
\end{tabular}]
{\haddockbegindoc
\(\mathcal{O}(1)\). Extract the first element of a list, which must be non-empty.\par
To disable the warning about partiality put \haddocktt{{\char '173}-{\char '43} OPTIONS{\char '137}GHC -Wno-x-partial -Wno-unrecognised-warning-flags {\char '43}-{\char '175}}
 at the top of the file. To disable it throughout a package put the same
 options into \haddocktt{ghc-options} section of Cabal file. To disable it in GHCi
 put \haddocktt{:set -Wno-x-partial -Wno-unrecognised-warning-flags} into \haddocktt{{\char '176}/.ghci} config file.
 See also the \href{https://github.com/haskell/core-libraries-committee/blob/main/guides/warning-for-head-and-tail.md}{migration guide}.\par
\subsubsection*{\textbf{Examples}}
\begin{quote}
{\haddockverb\begin{verbatim}
>>> head [1, 2, 3]
1

\end{verbatim}}
\end{quote}
\begin{quote}
{\haddockverb\begin{verbatim}
>>> head [1..]
1

\end{verbatim}}
\end{quote}
\begin{quote}
{\haddockverb\begin{verbatim}
>>> head []
*** Exception: Prelude.head: empty list

\end{verbatim}}
\end{quote}}
\end{haddockdesc}
\begin{haddockdesc}
\item[\begin{tabular}{@{}l}
last :: HasCallStack => {\char 91}a{\char 93} -> a
\end{tabular}]
{\haddockbegindoc
\(\mathcal{O}(n)\). Extract the last element of a list, which must be
 finite and non-empty.\par
WARNING: This function is partial. Consider using \haddockid{unsnoc} instead.\par
\subsubsection*{\textbf{Examples}}
\begin{quote}
{\haddockverb\begin{verbatim}
>>> last [1, 2, 3]
3

\end{verbatim}}
\end{quote}
\begin{quote}
{\haddockverb\begin{verbatim}
>>> last [1..]
* Hangs forever *

\end{verbatim}}
\end{quote}
\begin{quote}
{\haddockverb\begin{verbatim}
>>> last []
*** Exception: Prelude.last: empty list

\end{verbatim}}
\end{quote}}
\end{haddockdesc}
\begin{haddockdesc}
\item[\begin{tabular}{@{}l}
tail :: HasCallStack => {\char 91}a{\char 93} -> {\char 91}a{\char 93}
\end{tabular}]
{\haddockbegindoc
\(\mathcal{O}(1)\). Extract the elements after the head of a list, which
 must be non-empty.\par
To disable the warning about partiality put \haddocktt{{\char '173}-{\char '43} OPTIONS{\char '137}GHC -Wno-x-partial -Wno-unrecognised-warning-flags {\char '43}-{\char '175}}
 at the top of the file. To disable it throughout a package put the same
 options into \haddocktt{ghc-options} section of Cabal file. To disable it in GHCi
 put \haddocktt{:set -Wno-x-partial -Wno-unrecognised-warning-flags} into \haddocktt{{\char '176}/.ghci} config file.
 See also the \href{https://github.com/haskell/core-libraries-committee/blob/main/guides/warning-for-head-and-tail.md}{migration guide}.\par
\subsubsection*{\textbf{Examples}}
\begin{quote}
{\haddockverb\begin{verbatim}
>>> tail [1, 2, 3]
[2,3]

\end{verbatim}}
\end{quote}
\begin{quote}
{\haddockverb\begin{verbatim}
>>> tail [1]
[]

\end{verbatim}}
\end{quote}
\begin{quote}
{\haddockverb\begin{verbatim}
>>> tail []
*** Exception: Prelude.tail: empty list

\end{verbatim}}
\end{quote}}
\end{haddockdesc}
\begin{haddockdesc}
\item[\begin{tabular}{@{}l}
init :: HasCallStack => {\char 91}a{\char 93} -> {\char 91}a{\char 93}
\end{tabular}]
{\haddockbegindoc
\(\mathcal{O}(n)\). Return all the elements of a list except the last one.
 The list must be non-empty.\par
WARNING: This function is partial. Consider using \haddockid{unsnoc} instead.\par
\subsubsection*{\textbf{Examples}}
\begin{quote}
{\haddockverb\begin{verbatim}
>>> init [1, 2, 3]
[1,2]

\end{verbatim}}
\end{quote}
\begin{quote}
{\haddockverb\begin{verbatim}
>>> init [1]
[]

\end{verbatim}}
\end{quote}
\begin{quote}
{\haddockverb\begin{verbatim}
>>> init []
*** Exception: Prelude.init: empty list

\end{verbatim}}
\end{quote}}
\end{haddockdesc}
\begin{haddockdesc}
\item[\begin{tabular}{@{}l}
null :: Foldable t => t a -> Bool
\end{tabular}]
{\haddockbegindoc
Test whether the structure is empty.  The default implementation is
 Left-associative and lazy in both the initial element and the
 accumulator.  Thus optimised for structures where the first element can
 be accessed in constant time.  Structures where this is not the case
 should have a non-default implementation.\par
\subsubsection*{\textbf{Examples}}
Basic usage:\par
\begin{quote}
{\haddockverb\begin{verbatim}
>>> null []
True

\end{verbatim}}
\end{quote}
\begin{quote}
{\haddockverb\begin{verbatim}
>>> null [1]
False

\end{verbatim}}
\end{quote}
\haddockid{null} is expected to terminate even for infinite structures.
 The default implementation terminates provided the structure
 is bounded on the left (there is a leftmost element).\par
\begin{quote}
{\haddockverb\begin{verbatim}
>>> null [1..]
False

\end{verbatim}}
\end{quote}}
\end{haddockdesc}
\begin{haddockdesc}
\item[\begin{tabular}{@{}l}
length :: Foldable t => t a -> Int
\end{tabular}]
{\haddockbegindoc
Returns the size/length of a finite structure as an \haddockid{Int}.  The
 default implementation just counts elements starting with the leftmost.
 Instances for structures that can compute the element count faster
 than via element-by-element counting, should provide a specialised
 implementation.\par
\subsubsection*{\textbf{Examples}}
Basic usage:\par
\begin{quote}
{\haddockverb\begin{verbatim}
>>> length []
0

\end{verbatim}}
\end{quote}
\begin{quote}
{\haddockverb\begin{verbatim}
>>> length ['a', 'b', 'c']
3

>>> length [1..]
* Hangs forever *

\end{verbatim}}
\end{quote}}
\end{haddockdesc}
\section{List transformations}
\begin{haddockdesc}
\item[\begin{tabular}{@{}l}
map :: (a -> b) -> {\char 91}a{\char 93} -> {\char 91}b{\char 93}
\end{tabular}]
{\haddockbegindoc
\(\mathcal{O}(n)\). \haddockid{map} \haddocktt{f xs} is the list obtained by applying \haddocktt{f} to
 each element of \haddocktt{xs}, i.e.,\par
\begin{quote}
{\haddockverb\begin{verbatim}
map f [x1, x2, ..., xn] == [f x1, f x2, ..., f xn]
map f [x1, x2, ...] == [f x1, f x2, ...]\end{verbatim}}
\end{quote}
this means that \haddocktt{map id == id}\par
\subsubsection*{\textbf{Examples}}
\begin{quote}
{\haddockverb\begin{verbatim}
>>> map (+1) [1, 2, 3]
[2,3,4]

\end{verbatim}}
\end{quote}
\begin{quote}
{\haddockverb\begin{verbatim}
>>> map id [1, 2, 3]
[1,2,3]

\end{verbatim}}
\end{quote}
\begin{quote}
{\haddockverb\begin{verbatim}
>>> map (\n -> 3 * n + 1) [1, 2, 3]
[4,7,10]

\end{verbatim}}
\end{quote}}
\end{haddockdesc}
\begin{haddockdesc}
\item[\begin{tabular}{@{}l}
reverse :: {\char 91}a{\char 93} -> {\char 91}a{\char 93}
\end{tabular}]
{\haddockbegindoc
\(\mathcal{O}(n)\). \haddockid{reverse} \haddocktt{xs} returns the elements of \haddocktt{xs} in reverse order.
 \haddocktt{xs} must be finite.\par
\subsubsection*{\textbf{Laziness}}
\haddockid{reverse} is lazy in its elements.\par
\begin{quote}
{\haddockverb\begin{verbatim}
>>> head (reverse [undefined, 1])
1

\end{verbatim}}
\end{quote}
\begin{quote}
{\haddockverb\begin{verbatim}
>>> reverse (1 : 2 : undefined)
*** Exception: Prelude.undefined

\end{verbatim}}
\end{quote}
\subsubsection*{\textbf{Examples}}
\begin{quote}
{\haddockverb\begin{verbatim}
>>> reverse []
[]

\end{verbatim}}
\end{quote}
\begin{quote}
{\haddockverb\begin{verbatim}
>>> reverse [42]
[42]

\end{verbatim}}
\end{quote}
\begin{quote}
{\haddockverb\begin{verbatim}
>>> reverse [2,5,7]
[7,5,2]

\end{verbatim}}
\end{quote}
\begin{quote}
{\haddockverb\begin{verbatim}
>>> reverse [1..]
* Hangs forever *

\end{verbatim}}
\end{quote}}
\end{haddockdesc}
\begin{haddockdesc}
\item[\begin{tabular}{@{}l}
intersperse :: a -> {\char 91}a{\char 93} -> {\char 91}a{\char 93}
\end{tabular}]
{\haddockbegindoc
\(\mathcal{O}(n)\). The \haddockid{intersperse} function takes an element and a list
 and `intersperses' that element between the elements of the list.\par
\subsubsection*{\textbf{Laziness}}
\haddockid{intersperse} has the following properties\par
\begin{quote}
{\haddockverb\begin{verbatim}
>>> take 1 (intersperse undefined ('a' : undefined))
"a"

\end{verbatim}}
\end{quote}
\begin{quote}
{\haddockverb\begin{verbatim}
>>> take 2 (intersperse ',' ('a' : undefined))
"a*** Exception: Prelude.undefined

\end{verbatim}}
\end{quote}
\subsubsection*{\textbf{Examples}}
\begin{quote}
{\haddockverb\begin{verbatim}
>>> intersperse ',' "abcde"
"a,b,c,d,e"

\end{verbatim}}
\end{quote}
\begin{quote}
{\haddockverb\begin{verbatim}
>>> intersperse 1 [3, 4, 5]
[3,1,4,1,5]

\end{verbatim}}
\end{quote}}
\end{haddockdesc}
\begin{haddockdesc}
\item[\begin{tabular}{@{}l}
intercalate :: {\char 91}a{\char 93} -> {\char 91}{\char 91}a{\char 93}{\char 93} -> {\char 91}a{\char 93}
\end{tabular}]
{\haddockbegindoc
\haddockid{intercalate} \haddocktt{xs xss} is equivalent to \haddocktt{(\haddockid{concat} (\haddockid{intersperse} xs xss))}.
 It inserts the list \haddocktt{xs} in between the lists in \haddocktt{xss} and concatenates the
 result.\par
\subsubsection*{\textbf{Laziness}}
\haddockid{intercalate} has the following properties:\par
\begin{quote}
{\haddockverb\begin{verbatim}
>>> take 5 (intercalate undefined ("Lorem" : undefined))
"Lorem"

\end{verbatim}}
\end{quote}
\begin{quote}
{\haddockverb\begin{verbatim}
>>> take 6 (intercalate ", " ("Lorem" : undefined))
"Lorem*** Exception: Prelude.undefined

\end{verbatim}}
\end{quote}
\subsubsection*{\textbf{Examples}}
\begin{quote}
{\haddockverb\begin{verbatim}
>>> intercalate ", " ["Lorem", "ipsum", "dolor"]
"Lorem, ipsum, dolor"

\end{verbatim}}
\end{quote}
\begin{quote}
{\haddockverb\begin{verbatim}
>>> intercalate [0, 1] [[2, 3], [4, 5, 6], []]
[2,3,0,1,4,5,6,0,1]

\end{verbatim}}
\end{quote}
\begin{quote}
{\haddockverb\begin{verbatim}
>>> intercalate [1, 2, 3] [[], []]
[1,2,3]

\end{verbatim}}
\end{quote}}
\end{haddockdesc}
\begin{haddockdesc}
\item[\begin{tabular}{@{}l}
transpose :: {\char 91}{\char 91}a{\char 93}{\char 93} -> {\char 91}{\char 91}a{\char 93}{\char 93}
\end{tabular}]
{\haddockbegindoc
The \haddockid{transpose} function transposes the rows and columns of its argument.\par
\subsubsection*{\textbf{Laziness}}
\haddockid{transpose} is lazy in its elements\par
\begin{quote}
{\haddockverb\begin{verbatim}
>>> take 1 (transpose ['a' : undefined, 'b' : undefined])
["ab"]

\end{verbatim}}
\end{quote}
\subsubsection*{\textbf{Examples}}
\begin{quote}
{\haddockverb\begin{verbatim}
>>> transpose [[1,2,3],[4,5,6]]
[[1,4],[2,5],[3,6]]

\end{verbatim}}
\end{quote}
If some of the rows are shorter than the following rows, their elements are skipped:\par
\begin{quote}
{\haddockverb\begin{verbatim}
>>> transpose [[10,11],[20],[],[30,31,32]]
[[10,20,30],[11,31],[32]]

\end{verbatim}}
\end{quote}
For this reason the outer list must be finite; otherwise \haddockid{transpose} hangs:\par
\begin{quote}
{\haddockverb\begin{verbatim}
>>> transpose (repeat [])
* Hangs forever *

\end{verbatim}}
\end{quote}}
\end{haddockdesc}
\begin{haddockdesc}
\item[\begin{tabular}{@{}l}
subsequences :: {\char 91}a{\char 93} -> {\char 91}{\char 91}a{\char 93}{\char 93}
\end{tabular}]
{\haddockbegindoc
The \haddockid{subsequences} function returns the list of all subsequences of the argument.\par
\subsubsection*{\textbf{Laziness}}
\haddockid{subsequences} does not look ahead unless it must:\par
\begin{quote}
{\haddockverb\begin{verbatim}
>>> take 1 (subsequences undefined)
[[]]

>>> take 2 (subsequences ('a' : undefined))
["","a"]

\end{verbatim}}
\end{quote}
\subsubsection*{\textbf{Examples}}
\begin{quote}
{\haddockverb\begin{verbatim}
>>> subsequences "abc"
["","a","b","ab","c","ac","bc","abc"]

\end{verbatim}}
\end{quote}
This function is productive on infinite inputs:\par
\begin{quote}
{\haddockverb\begin{verbatim}
>>> take 8 $ subsequences ['a'..]
["","a","b","ab","c","ac","bc","abc"]

\end{verbatim}}
\end{quote}}
\end{haddockdesc}
\begin{haddockdesc}
\item[\begin{tabular}{@{}l}
permutations :: {\char 91}a{\char 93} -> {\char 91}{\char 91}a{\char 93}{\char 93}
\end{tabular}]
{\haddockbegindoc
The \haddockid{permutations} function returns the list of all permutations of the argument.\par
Note that the order of permutations is not lexicographic.
 It satisfies the following property:\par
\begin{quote}
{\haddockverb\begin{verbatim}
map (take n) (take (product [1..n]) (permutations ([1..n] ++ undefined))) == permutations [1..n]\end{verbatim}}
\end{quote}
\subsubsection*{\textbf{Laziness}}
The \haddockid{permutations} function is maximally lazy:
 for each \haddocktt{n}, the value of \haddocktt{\haddockid{permutations} xs} starts with those permutations
 that permute \haddocktt{\haddockid{take} n xs} and keep \haddocktt{\haddockid{drop} n xs}.\par
\subsubsection*{\textbf{Examples}}
\begin{quote}
{\haddockverb\begin{verbatim}
>>> permutations "abc"
["abc","bac","cba","bca","cab","acb"]

\end{verbatim}}
\end{quote}
\begin{quote}
{\haddockverb\begin{verbatim}
>>> permutations [1, 2]
[[1,2],[2,1]]

\end{verbatim}}
\end{quote}
\begin{quote}
{\haddockverb\begin{verbatim}
>>> permutations []
[[]]

\end{verbatim}}
\end{quote}
This function is productive on infinite inputs:\par
\begin{quote}
{\haddockverb\begin{verbatim}
>>> take 6 $ map (take 3) $ permutations ['a'..]
["abc","bac","cba","bca","cab","acb"]

\end{verbatim}}
\end{quote}}
\end{haddockdesc}
\section{Reducing lists (folds)}
\begin{haddockdesc}
\item[\begin{tabular}{@{}l}
foldl :: Foldable t => (b -> a -> b) -> b -> t a -> b
\end{tabular}]
{\haddockbegindoc
Left-associative fold of a structure, lazy in the accumulator.  This
 is rarely what you want, but can work well for structures with efficient
 right-to-left sequencing and an operator that is lazy in its left
 argument.\par
In the case of lists, \haddockid{foldl}, when applied to a binary operator, a
 starting value (typically the left-identity of the operator), and a
 list, reduces the list using the binary operator, from left to right:\par
\begin{quote}
{\haddockverb\begin{verbatim}
foldl f z [x1, x2, ..., xn] == (...((z `f` x1) `f` x2) `f`...) `f` xn\end{verbatim}}
\end{quote}
Note that to produce the outermost application of the operator the
 entire input list must be traversed.  Like all left-associative folds,
 \haddockid{foldl} will diverge if given an infinite list.\par
If you want an efficient strict left-fold, you probably want to use
 \haddockid{foldl'} instead of \haddockid{foldl}.  The reason for this is that the latter
 does not force the \emph{inner} results (e.g. \haddocktt{z `f` x1} in the above
 example) before applying them to the operator (e.g. to \haddocktt{(`f` x2)}).
 This results in a thunk chain \emph{O(n)} elements long, which then must be
 evaluated from the outside-in.\par
For a general \haddockid{Foldable} structure this should be semantically identical
 to:\par
\begin{quote}
{\haddockverb\begin{verbatim}
foldl f z = foldl f z . toList\end{verbatim}}
\end{quote}
\subsubsection*{\textbf{Examples}}
The first example is a strict fold, which in practice is best performed
 with \haddockid{foldl'}.\par
\begin{quote}
{\haddockverb\begin{verbatim}
>>> foldl (+) 42 [1,2,3,4]
52

\end{verbatim}}
\end{quote}
Though the result below is lazy, the input is reversed before prepending
 it to the initial accumulator, so corecursion begins only after traversing
 the entire input string.\par
\begin{quote}
{\haddockverb\begin{verbatim}
>>> foldl (\acc c -> c : acc) "abcd" "efgh"
"hgfeabcd"

\end{verbatim}}
\end{quote}
A left fold of a structure that is infinite on the right cannot
 terminate, even when for any finite input the fold just returns the
 initial accumulator:\par
\begin{quote}
{\haddockverb\begin{verbatim}
>>> foldl (\a _ -> a) 0 $ repeat 1
* Hangs forever *

\end{verbatim}}
\end{quote}
WARNING: When it comes to lists, you always want to use either \haddockid{foldl'} or \haddockid{foldr} instead.\par}
\end{haddockdesc}
\begin{haddockdesc}
\item[\begin{tabular}{@{}l}
foldl' :: Foldable t => (b -> a -> b) -> b -> t a -> b
\end{tabular}]
{\haddockbegindoc
Left-associative fold of a structure but with strict application of
 the operator.\par
This ensures that each step of the fold is forced to Weak Head Normal
 Form before being applied, avoiding the collection of thunks that would
 otherwise occur.  This is often what you want to strictly reduce a
 finite structure to a single strict result (e.g. \haddockid{sum}).\par
For a general \haddockid{Foldable} structure this should be semantically identical
 to,\par
\begin{quote}
{\haddockverb\begin{verbatim}
foldl' f z = foldl' f z . toList\end{verbatim}}
\end{quote}}
\end{haddockdesc}
\begin{haddockdesc}
\item[\begin{tabular}{@{}l}
foldl1 :: Foldable t => (a -> a -> a) -> t a -> a
\end{tabular}]
{\haddockbegindoc
A variant of \haddockid{foldl} that has no base case,
 and thus may only be applied to non-empty structures.\par
This function is non-total and will raise a runtime exception if the
 structure happens to be empty.\par
\begin{quote}
{\haddockverb\begin{verbatim}
foldl1 f = foldl1 f . toList\end{verbatim}}
\end{quote}
\subsubsection*{\textbf{Examples}}
Basic usage:\par
\begin{quote}
{\haddockverb\begin{verbatim}
>>> foldl1 (+) [1..4]
10

\end{verbatim}}
\end{quote}
\begin{quote}
{\haddockverb\begin{verbatim}
>>> foldl1 (+) []
*** Exception: Prelude.foldl1: empty list

\end{verbatim}}
\end{quote}
\begin{quote}
{\haddockverb\begin{verbatim}
>>> foldl1 (+) Nothing
*** Exception: foldl1: empty structure

\end{verbatim}}
\end{quote}
\begin{quote}
{\haddockverb\begin{verbatim}
>>> foldl1 (-) [1..4]
-8

\end{verbatim}}
\end{quote}
\begin{quote}
{\haddockverb\begin{verbatim}
>>> foldl1 (&&) [True, False, True, True]
False

\end{verbatim}}
\end{quote}
\begin{quote}
{\haddockverb\begin{verbatim}
>>> foldl1 (||) [False, False, True, True]
True

\end{verbatim}}
\end{quote}
\begin{quote}
{\haddockverb\begin{verbatim}
>>> foldl1 (+) [1..]
* Hangs forever *

\end{verbatim}}
\end{quote}}
\end{haddockdesc}
\begin{haddockdesc}
\item[\begin{tabular}{@{}l}
foldl1' :: HasCallStack => (a -> a -> a) -> {\char 91}a{\char 93} -> a
\end{tabular}]
{\haddockbegindoc
A strict version of \haddockid{foldl1}.\par}
\end{haddockdesc}
\begin{haddockdesc}
\item[\begin{tabular}{@{}l}
foldr :: Foldable t => (a -> b -> b) -> b -> t a -> b
\end{tabular}]
{\haddockbegindoc
Right-associative fold of a structure, lazy in the accumulator.\par
In the case of lists, \haddockid{foldr}, when applied to a binary operator, a
 starting value (typically the right-identity of the operator), and a
 list, reduces the list using the binary operator, from right to left:\par
\begin{quote}
{\haddockverb\begin{verbatim}
foldr f z [x1, x2, ..., xn] == x1 `f` (x2 `f` ... (xn `f` z)...)\end{verbatim}}
\end{quote}
Note that since the head of the resulting expression is produced by an
 application of the operator to the first element of the list, given an
 operator lazy in its right argument, \haddockid{foldr} can produce a terminating
 expression from an unbounded list.\par
For a general \haddockid{Foldable} structure this should be semantically identical
 to,\par
\begin{quote}
{\haddockverb\begin{verbatim}
foldr f z = foldr f z . toList\end{verbatim}}
\end{quote}
\subsubsection*{\textbf{Examples}}
Basic usage:\par
\begin{quote}
{\haddockverb\begin{verbatim}
>>> foldr (||) False [False, True, False]
True

\end{verbatim}}
\end{quote}
\begin{quote}
{\haddockverb\begin{verbatim}
>>> foldr (||) False []
False

\end{verbatim}}
\end{quote}
\begin{quote}
{\haddockverb\begin{verbatim}
>>> foldr (\c acc -> acc ++ [c]) "foo" ['a', 'b', 'c', 'd']
"foodcba"

\end{verbatim}}
\end{quote}
\subsubsection*{Infinite structures}
⚠️ Applying \haddockid{foldr} to infinite structures usually doesn't terminate.\par
It may still terminate under one of the following conditions:\par
\vbox{\begin{itemize}
\item
the folding function is short-circuiting\par
\item
the folding function is lazy on its second argument\par
\end{itemize}}
\subsubsection*{Short-circuiting}
\haddocktt{(\haddockid{||})} short-circuits on \haddockid{True} values, so the following terminates
 because there is a \haddockid{True} value finitely far from the left side:\par
\begin{quote}
{\haddockverb\begin{verbatim}
>>> foldr (||) False (True : repeat False)
True

\end{verbatim}}
\end{quote}
But the following doesn't terminate:\par
\begin{quote}
{\haddockverb\begin{verbatim}
>>> foldr (||) False (repeat False ++ [True])
* Hangs forever *

\end{verbatim}}
\end{quote}
\subsubsection*{Laziness in the second argument}
Applying \haddockid{foldr} to infinite structures terminates when the operator is
 lazy in its second argument (the initial accumulator is never used in
 this case, and so could be left \haddockid{undefined}, but \haddocktt{{\char 91}{\char 93}} is more clear):\par
\begin{quote}
{\haddockverb\begin{verbatim}
>>> take 5 $ foldr (\i acc -> i : fmap (+3) acc) [] (repeat 1)
[1,4,7,10,13]

\end{verbatim}}
\end{quote}}
\end{haddockdesc}
\begin{haddockdesc}
\item[\begin{tabular}{@{}l}
foldr1 :: Foldable t => (a -> a -> a) -> t a -> a
\end{tabular}]
{\haddockbegindoc
A variant of \haddockid{foldr} that has no base case,
 and thus may only be applied to non-empty structures.\par
This function is non-total and will raise a runtime exception if the
 structure happens to be empty.\par
\subsubsection*{\textbf{Examples}}
Basic usage:\par
\begin{quote}
{\haddockverb\begin{verbatim}
>>> foldr1 (+) [1..4]
10

\end{verbatim}}
\end{quote}
\begin{quote}
{\haddockverb\begin{verbatim}
>>> foldr1 (+) []
Exception: Prelude.foldr1: empty list

\end{verbatim}}
\end{quote}
\begin{quote}
{\haddockverb\begin{verbatim}
>>> foldr1 (+) Nothing
*** Exception: foldr1: empty structure

\end{verbatim}}
\end{quote}
\begin{quote}
{\haddockverb\begin{verbatim}
>>> foldr1 (-) [1..4]
-2

\end{verbatim}}
\end{quote}
\begin{quote}
{\haddockverb\begin{verbatim}
>>> foldr1 (&&) [True, False, True, True]
False

\end{verbatim}}
\end{quote}
\begin{quote}
{\haddockverb\begin{verbatim}
>>> foldr1 (||) [False, False, True, True]
True

\end{verbatim}}
\end{quote}
\begin{quote}
{\haddockverb\begin{verbatim}
>>> foldr1 (+) [1..]
* Hangs forever *

\end{verbatim}}
\end{quote}}
\end{haddockdesc}
\section{Special folds}
\begin{haddockdesc}
\item[\begin{tabular}{@{}l}
concat :: Foldable t => t {\char 91}a{\char 93} -> {\char 91}a{\char 93}
\end{tabular}]
{\haddockbegindoc
The concatenation of all the elements of a container of lists.\par
\subsubsection*{\textbf{Examples}}
Basic usage:\par
\begin{quote}
{\haddockverb\begin{verbatim}
>>> concat (Just [1, 2, 3])
[1,2,3]

\end{verbatim}}
\end{quote}
\begin{quote}
{\haddockverb\begin{verbatim}
>>> concat (Left 42)
[]

\end{verbatim}}
\end{quote}
\begin{quote}
{\haddockverb\begin{verbatim}
>>> concat [[1, 2, 3], [4, 5], [6], []]
[1,2,3,4,5,6]

\end{verbatim}}
\end{quote}}
\end{haddockdesc}
\begin{haddockdesc}
\item[\begin{tabular}{@{}l}
concatMap :: Foldable t => (a -> {\char 91}b{\char 93}) -> t a -> {\char 91}b{\char 93}
\end{tabular}]
{\haddockbegindoc
Map a function over all the elements of a container and concatenate
 the resulting lists.\par
\subsubsection*{\textbf{Examples}}
Basic usage:\par
\begin{quote}
{\haddockverb\begin{verbatim}
>>> concatMap (take 3) [[1..], [10..], [100..], [1000..]]
[1,2,3,10,11,12,100,101,102,1000,1001,1002]

\end{verbatim}}
\end{quote}
\begin{quote}
{\haddockverb\begin{verbatim}
>>> concatMap (take 3) (Just [1..])
[1,2,3]

\end{verbatim}}
\end{quote}}
\end{haddockdesc}
\begin{haddockdesc}
\item[\begin{tabular}{@{}l}
and :: Foldable t => t Bool -> Bool
\end{tabular}]
{\haddockbegindoc
\haddockid{and} returns the conjunction of a container of Bools.  For the
 result to be \haddockid{True}, the container must be finite; \haddockid{False}, however,
 results from a \haddockid{False} value finitely far from the left end.\par
\subsubsection*{\textbf{Examples}}
Basic usage:\par
\begin{quote}
{\haddockverb\begin{verbatim}
>>> and []
True

\end{verbatim}}
\end{quote}
\begin{quote}
{\haddockverb\begin{verbatim}
>>> and [True]
True

\end{verbatim}}
\end{quote}
\begin{quote}
{\haddockverb\begin{verbatim}
>>> and [False]
False

\end{verbatim}}
\end{quote}
\begin{quote}
{\haddockverb\begin{verbatim}
>>> and [True, True, False]
False

\end{verbatim}}
\end{quote}
\begin{quote}
{\haddockverb\begin{verbatim}
>>> and (False : repeat True) -- Infinite list [False,True,True,True,...
False

\end{verbatim}}
\end{quote}
\begin{quote}
{\haddockverb\begin{verbatim}
>>> and (repeat True)
* Hangs forever *

\end{verbatim}}
\end{quote}}
\end{haddockdesc}
\begin{haddockdesc}
\item[\begin{tabular}{@{}l}
or :: Foldable t => t Bool -> Bool
\end{tabular}]
{\haddockbegindoc
\haddockid{or} returns the disjunction of a container of Bools.  For the
 result to be \haddockid{False}, the container must be finite; \haddockid{True}, however,
 results from a \haddockid{True} value finitely far from the left end.\par
\subsubsection*{\textbf{Examples}}
Basic usage:\par
\begin{quote}
{\haddockverb\begin{verbatim}
>>> or []
False

\end{verbatim}}
\end{quote}
\begin{quote}
{\haddockverb\begin{verbatim}
>>> or [True]
True

\end{verbatim}}
\end{quote}
\begin{quote}
{\haddockverb\begin{verbatim}
>>> or [False]
False

\end{verbatim}}
\end{quote}
\begin{quote}
{\haddockverb\begin{verbatim}
>>> or [True, True, False]
True

\end{verbatim}}
\end{quote}
\begin{quote}
{\haddockverb\begin{verbatim}
>>> or (True : repeat False) -- Infinite list [True,False,False,False,...
True

\end{verbatim}}
\end{quote}
\begin{quote}
{\haddockverb\begin{verbatim}
>>> or (repeat False)
* Hangs forever *

\end{verbatim}}
\end{quote}}
\end{haddockdesc}
\begin{haddockdesc}
\item[\begin{tabular}{@{}l}
any :: Foldable t => (a -> Bool) -> t a -> Bool
\end{tabular}]
{\haddockbegindoc
Determines whether any element of the structure satisfies the predicate.\par
\subsubsection*{\textbf{Examples}}
Basic usage:\par
\begin{quote}
{\haddockverb\begin{verbatim}
>>> any (> 3) []
False

\end{verbatim}}
\end{quote}
\begin{quote}
{\haddockverb\begin{verbatim}
>>> any (> 3) [1,2]
False

\end{verbatim}}
\end{quote}
\begin{quote}
{\haddockverb\begin{verbatim}
>>> any (> 3) [1,2,3,4,5]
True

\end{verbatim}}
\end{quote}
\begin{quote}
{\haddockverb\begin{verbatim}
>>> any (> 3) [1..]
True

\end{verbatim}}
\end{quote}
\begin{quote}
{\haddockverb\begin{verbatim}
>>> any (> 3) [0, -1..]
* Hangs forever *

\end{verbatim}}
\end{quote}}
\end{haddockdesc}
\begin{haddockdesc}
\item[\begin{tabular}{@{}l}
all :: Foldable t => (a -> Bool) -> t a -> Bool
\end{tabular}]
{\haddockbegindoc
Determines whether all elements of the structure satisfy the predicate.\par
\subsubsection*{\textbf{Examples}}
Basic usage:\par
\begin{quote}
{\haddockverb\begin{verbatim}
>>> all (> 3) []
True

\end{verbatim}}
\end{quote}
\begin{quote}
{\haddockverb\begin{verbatim}
>>> all (> 3) [1,2]
False

\end{verbatim}}
\end{quote}
\begin{quote}
{\haddockverb\begin{verbatim}
>>> all (> 3) [1,2,3,4,5]
False

\end{verbatim}}
\end{quote}
\begin{quote}
{\haddockverb\begin{verbatim}
>>> all (> 3) [1..]
False

\end{verbatim}}
\end{quote}
\begin{quote}
{\haddockverb\begin{verbatim}
>>> all (> 3) [4..]
* Hangs forever *

\end{verbatim}}
\end{quote}}
\end{haddockdesc}
\begin{haddockdesc}
\item[\begin{tabular}{@{}l}
sum :: (Foldable t, Num a) => t a -> a
\end{tabular}]
{\haddockbegindoc
The \haddockid{sum} function computes the sum of the numbers of a structure.\par
\subsubsection*{\textbf{Examples}}
Basic usage:\par
\begin{quote}
{\haddockverb\begin{verbatim}
>>> sum []
0

\end{verbatim}}
\end{quote}
\begin{quote}
{\haddockverb\begin{verbatim}
>>> sum [42]
42

\end{verbatim}}
\end{quote}
\begin{quote}
{\haddockverb\begin{verbatim}
>>> sum [1..10]
55

\end{verbatim}}
\end{quote}
\begin{quote}
{\haddockverb\begin{verbatim}
>>> sum [4.1, 2.0, 1.7]
7.8

\end{verbatim}}
\end{quote}
\begin{quote}
{\haddockverb\begin{verbatim}
>>> sum [1..]
* Hangs forever *

\end{verbatim}}
\end{quote}}
\end{haddockdesc}
\begin{haddockdesc}
\item[\begin{tabular}{@{}l}
product :: (Foldable t, Num a) => t a -> a
\end{tabular}]
{\haddockbegindoc
The \haddockid{product} function computes the product of the numbers of a
 structure.\par
\subsubsection*{\textbf{Examples}}
Basic usage:\par
\begin{quote}
{\haddockverb\begin{verbatim}
>>> product []
1

\end{verbatim}}
\end{quote}
\begin{quote}
{\haddockverb\begin{verbatim}
>>> product [42]
42

\end{verbatim}}
\end{quote}
\begin{quote}
{\haddockverb\begin{verbatim}
>>> product [1..10]
3628800

\end{verbatim}}
\end{quote}
\begin{quote}
{\haddockverb\begin{verbatim}
>>> product [4.1, 2.0, 1.7]
13.939999999999998

\end{verbatim}}
\end{quote}
\begin{quote}
{\haddockverb\begin{verbatim}
>>> product [1..]
* Hangs forever *

\end{verbatim}}
\end{quote}}
\end{haddockdesc}
\begin{haddockdesc}
\item[\begin{tabular}{@{}l}
maximum :: (Foldable t, Ord a) => t a -> a
\end{tabular}]
{\haddockbegindoc
The largest element of a non-empty structure.\par
This function is non-total and will raise a runtime exception if the
 structure happens to be empty.  A structure that supports random access
 and maintains its elements in order should provide a specialised
 implementation to return the maximum in faster than linear time.\par
\subsubsection*{\textbf{Examples}}
Basic usage:\par
\begin{quote}
{\haddockverb\begin{verbatim}
>>> maximum [1..10]
10

\end{verbatim}}
\end{quote}
\begin{quote}
{\haddockverb\begin{verbatim}
>>> maximum []
*** Exception: Prelude.maximum: empty list

\end{verbatim}}
\end{quote}
\begin{quote}
{\haddockverb\begin{verbatim}
>>> maximum Nothing
*** Exception: maximum: empty structure

\end{verbatim}}
\end{quote}
WARNING: This function is partial for possibly-empty structures like lists.\par}
\end{haddockdesc}
\begin{haddockdesc}
\item[\begin{tabular}{@{}l}
minimum :: (Foldable t, Ord a) => t a -> a
\end{tabular}]
{\haddockbegindoc
The least element of a non-empty structure.\par
This function is non-total and will raise a runtime exception if the
 structure happens to be empty.  A structure that supports random access
 and maintains its elements in order should provide a specialised
 implementation to return the minimum in faster than linear time.\par
\subsubsection*{\textbf{Examples}}
Basic usage:\par
\begin{quote}
{\haddockverb\begin{verbatim}
>>> minimum [1..10]
1

\end{verbatim}}
\end{quote}
\begin{quote}
{\haddockverb\begin{verbatim}
>>> minimum []
*** Exception: Prelude.minimum: empty list

\end{verbatim}}
\end{quote}
\begin{quote}
{\haddockverb\begin{verbatim}
>>> minimum Nothing
*** Exception: minimum: empty structure

\end{verbatim}}
\end{quote}
WARNING: This function is partial for possibly-empty structures like lists.\par}
\end{haddockdesc}
\section{Building lists}
\subsection{Scans}
\begin{haddockdesc}
\item[\begin{tabular}{@{}l}
scanl :: (b -> a -> b) -> b -> {\char 91}a{\char 93} -> {\char 91}b{\char 93}
\end{tabular}]
{\haddockbegindoc
\(\mathcal{O}(n)\). \haddockid{scanl} is similar to \haddockid{foldl}, but returns a list of
 successive reduced values from the left:\par
\begin{quote}
{\haddockverb\begin{verbatim}
scanl f z [x1, x2, ...] == [z, z `f` x1, (z `f` x1) `f` x2, ...]\end{verbatim}}
\end{quote}
Note that\par
\begin{quote}
{\haddockverb\begin{verbatim}
last (scanl f z xs) == foldl f z xs\end{verbatim}}
\end{quote}
\subsubsection*{\textbf{Examples}}
\begin{quote}
{\haddockverb\begin{verbatim}
>>> scanl (+) 0 [1..4]
[0,1,3,6,10]

\end{verbatim}}
\end{quote}
\begin{quote}
{\haddockverb\begin{verbatim}
>>> scanl (+) 42 []
[42]

\end{verbatim}}
\end{quote}
\begin{quote}
{\haddockverb\begin{verbatim}
>>> scanl (-) 100 [1..4]
[100,99,97,94,90]

\end{verbatim}}
\end{quote}
\begin{quote}
{\haddockverb\begin{verbatim}
>>> scanl (\reversedString nextChar -> nextChar : reversedString) "foo" ['a', 'b', 'c', 'd']
["foo","afoo","bafoo","cbafoo","dcbafoo"]

\end{verbatim}}
\end{quote}
\begin{quote}
{\haddockverb\begin{verbatim}
>>> take 10 (scanl (+) 0 [1..])
[0,1,3,6,10,15,21,28,36,45]

\end{verbatim}}
\end{quote}
\begin{quote}
{\haddockverb\begin{verbatim}
>>> take 1 (scanl undefined 'a' undefined)
"a"

\end{verbatim}}
\end{quote}}
\end{haddockdesc}
\begin{haddockdesc}
\item[\begin{tabular}{@{}l}
scanl1 :: (a -> a -> a) -> {\char 91}a{\char 93} -> {\char 91}a{\char 93}
\end{tabular}]
{\haddockbegindoc
\(\mathcal{O}(n)\). \haddockid{scanl1} is a variant of \haddockid{scanl} that has no starting
 value argument:\par
\begin{quote}
{\haddockverb\begin{verbatim}
scanl1 f [x1, x2, ...] == [x1, x1 `f` x2, ...]\end{verbatim}}
\end{quote}
\subsubsection*{\textbf{Examples}}
\begin{quote}
{\haddockverb\begin{verbatim}
>>> scanl1 (+) [1..4]
[1,3,6,10]

\end{verbatim}}
\end{quote}
\begin{quote}
{\haddockverb\begin{verbatim}
>>> scanl1 (+) []
[]

\end{verbatim}}
\end{quote}
\begin{quote}
{\haddockverb\begin{verbatim}
>>> scanl1 (-) [1..4]
[1,-1,-4,-8]

\end{verbatim}}
\end{quote}
\begin{quote}
{\haddockverb\begin{verbatim}
>>> scanl1 (&&) [True, False, True, True]
[True,False,False,False]

\end{verbatim}}
\end{quote}
\begin{quote}
{\haddockverb\begin{verbatim}
>>> scanl1 (||) [False, False, True, True]
[False,False,True,True]

\end{verbatim}}
\end{quote}
\begin{quote}
{\haddockverb\begin{verbatim}
>>> take 10 (scanl1 (+) [1..])
[1,3,6,10,15,21,28,36,45,55]

\end{verbatim}}
\end{quote}
\begin{quote}
{\haddockverb\begin{verbatim}
>>> take 1 (scanl1 undefined ('a' : undefined))
"a"

\end{verbatim}}
\end{quote}}
\end{haddockdesc}
\begin{haddockdesc}
\item[\begin{tabular}{@{}l}
scanr :: (a -> b -> b) -> b -> {\char 91}a{\char 93} -> {\char 91}b{\char 93}
\end{tabular}]
{\haddockbegindoc
\(\mathcal{O}(n)\). \haddockid{scanr} is the right-to-left dual of \haddockid{scanl}. Note that the order of parameters on the accumulating function are reversed compared to \haddockid{scanl}.
 Also note that\par
\begin{quote}
{\haddockverb\begin{verbatim}
head (scanr f z xs) == foldr f z xs.\end{verbatim}}
\end{quote}
\subsubsection*{\textbf{Examples}}
\begin{quote}
{\haddockverb\begin{verbatim}
>>> scanr (+) 0 [1..4]
[10,9,7,4,0]

\end{verbatim}}
\end{quote}
\begin{quote}
{\haddockverb\begin{verbatim}
>>> scanr (+) 42 []
[42]

\end{verbatim}}
\end{quote}
\begin{quote}
{\haddockverb\begin{verbatim}
>>> scanr (-) 100 [1..4]
[98,-97,99,-96,100]

\end{verbatim}}
\end{quote}
\begin{quote}
{\haddockverb\begin{verbatim}
>>> scanr (\nextChar reversedString -> nextChar : reversedString) "foo" ['a', 'b', 'c', 'd']
["abcdfoo","bcdfoo","cdfoo","dfoo","foo"]

\end{verbatim}}
\end{quote}
\begin{quote}
{\haddockverb\begin{verbatim}
>>> force $ scanr (+) 0 [1..]
*** Exception: stack overflow

\end{verbatim}}
\end{quote}}
\end{haddockdesc}
\begin{haddockdesc}
\item[\begin{tabular}{@{}l}
scanr1 :: (a -> a -> a) -> {\char 91}a{\char 93} -> {\char 91}a{\char 93}
\end{tabular}]
{\haddockbegindoc
\(\mathcal{O}(n)\). \haddockid{scanr1} is a variant of \haddockid{scanr} that has no starting
 value argument.\par
\subsubsection*{\textbf{Examples}}
\begin{quote}
{\haddockverb\begin{verbatim}
>>> scanr1 (+) [1..4]
[10,9,7,4]

\end{verbatim}}
\end{quote}
\begin{quote}
{\haddockverb\begin{verbatim}
>>> scanr1 (+) []
[]

\end{verbatim}}
\end{quote}
\begin{quote}
{\haddockverb\begin{verbatim}
>>> scanr1 (-) [1..4]
[-2,3,-1,4]

\end{verbatim}}
\end{quote}
\begin{quote}
{\haddockverb\begin{verbatim}
>>> scanr1 (&&) [True, False, True, True]
[False,False,True,True]

\end{verbatim}}
\end{quote}
\begin{quote}
{\haddockverb\begin{verbatim}
>>> scanr1 (||) [True, True, False, False]
[True,True,False,False]

\end{verbatim}}
\end{quote}
\begin{quote}
{\haddockverb\begin{verbatim}
>>> force $ scanr1 (+) [1..]
*** Exception: stack overflow

\end{verbatim}}
\end{quote}}
\end{haddockdesc}
\subsection{Accumulating maps}
\begin{haddockdesc}
\item[\begin{tabular}{@{}l}
mapAccumL :: Traversable t => (s -> a -> (s, b)) -> s -> t a -> (s, t b)
\end{tabular}]
{\haddockbegindoc
The \haddockid{mapAccumL} function behaves like a combination of \haddockid{fmap}
 and \haddockid{foldl}; it applies a function to each element of a structure,
 passing an accumulating parameter from left to right, and returning
 a final value of this accumulator together with the new structure.\par
\subsubsection*{\textbf{Examples}}
Basic usage:\par
\begin{quote}
{\haddockverb\begin{verbatim}
>>> mapAccumL (\a b -> (a + b, a)) 0 [1..10]
(55,[0,1,3,6,10,15,21,28,36,45])

\end{verbatim}}
\end{quote}
\begin{quote}
{\haddockverb\begin{verbatim}
>>> mapAccumL (\a b -> (a <> show b, a)) "0" [1..5]
("012345",["0","01","012","0123","01234"])

\end{verbatim}}
\end{quote}}
\end{haddockdesc}
\begin{haddockdesc}
\item[\begin{tabular}{@{}l}
mapAccumR :: Traversable t => (s -> a -> (s, b)) -> s -> t a -> (s, t b)
\end{tabular}]
{\haddockbegindoc
The \haddockid{mapAccumR} function behaves like a combination of \haddockid{fmap}
 and \haddockid{foldr}; it applies a function to each element of a structure,
 passing an accumulating parameter from right to left, and returning
 a final value of this accumulator together with the new structure.\par
\subsubsection*{\textbf{Examples}}
Basic usage:\par
\begin{quote}
{\haddockverb\begin{verbatim}
>>> mapAccumR (\a b -> (a + b, a)) 0 [1..10]
(55,[54,52,49,45,40,34,27,19,10,0])

\end{verbatim}}
\end{quote}
\begin{quote}
{\haddockverb\begin{verbatim}
>>> mapAccumR (\a b -> (a <> show b, a)) "0" [1..5]
("054321",["05432","0543","054","05","0"])

\end{verbatim}}
\end{quote}}
\end{haddockdesc}
\subsection{Infinite lists}
\begin{haddockdesc}
\item[\begin{tabular}{@{}l}
iterate :: (a -> a) -> a -> {\char 91}a{\char 93}
\end{tabular}]
{\haddockbegindoc
\haddockid{iterate} \haddocktt{f x} returns an infinite list of repeated applications
 of \haddocktt{f} to \haddocktt{x}:\par
\begin{quote}
{\haddockverb\begin{verbatim}
iterate f x == [x, f x, f (f x), ...]\end{verbatim}}
\end{quote}
\subsubsection*{\textbf{Laziness}}
Note that \haddockid{iterate} is lazy, potentially leading to thunk build-up if
 the consumer doesn't force each iterate. See \haddockid{iterate'} for a strict
 variant of this function.\par
\begin{quote}
{\haddockverb\begin{verbatim}
>>> take 1 $ iterate undefined 42
[42]

\end{verbatim}}
\end{quote}
\subsubsection*{\textbf{Examples}}
\begin{quote}
{\haddockverb\begin{verbatim}
>>> take 10 $ iterate not True
[True,False,True,False,True,False,True,False,True,False]

\end{verbatim}}
\end{quote}
\begin{quote}
{\haddockverb\begin{verbatim}
>>> take 10 $ iterate (+3) 42
[42,45,48,51,54,57,60,63,66,69]

\end{verbatim}}
\end{quote}
\haddocktt{iterate id == \haddockid{repeat}}:\par
\begin{quote}
{\haddockverb\begin{verbatim}
>>> take 10 $ iterate id 1
[1,1,1,1,1,1,1,1,1,1]

\end{verbatim}}
\end{quote}}
\end{haddockdesc}
\begin{haddockdesc}
\item[\begin{tabular}{@{}l}
repeat :: a -> {\char 91}a{\char 93}
\end{tabular}]
{\haddockbegindoc
\haddockid{repeat} \haddocktt{x} is an infinite list, with \haddocktt{x} the value of every element.\par
\subsubsection*{\textbf{Examples}}
\begin{quote}
{\haddockverb\begin{verbatim}
>>> take 10 $ repeat 17
[17,17,17,17,17,17,17,17,17, 17]

\end{verbatim}}
\end{quote}
\begin{quote}
{\haddockverb\begin{verbatim}
>>> repeat undefined
[*** Exception: Prelude.undefined

\end{verbatim}}
\end{quote}}
\end{haddockdesc}
\begin{haddockdesc}
\item[\begin{tabular}{@{}l}
replicate :: Int -> a -> {\char 91}a{\char 93}
\end{tabular}]
{\haddockbegindoc
\haddockid{replicate} \haddocktt{n x} is a list of length \haddocktt{n} with \haddocktt{x} the value of
 every element.
 It is an instance of the more general \haddockid{genericReplicate},
 in which \haddocktt{n} may be of any integral type.\par
\subsubsection*{\textbf{Examples}}
\begin{quote}
{\haddockverb\begin{verbatim}
>>> replicate 0 True
[]

\end{verbatim}}
\end{quote}
\begin{quote}
{\haddockverb\begin{verbatim}
>>> replicate (-1) True
[]

\end{verbatim}}
\end{quote}
\begin{quote}
{\haddockverb\begin{verbatim}
>>> replicate 4 True
[True,True,True,True]

\end{verbatim}}
\end{quote}}
\end{haddockdesc}
\begin{haddockdesc}
\item[\begin{tabular}{@{}l}
cycle :: HasCallStack => {\char 91}a{\char 93} -> {\char 91}a{\char 93}
\end{tabular}]
{\haddockbegindoc
\haddockid{cycle} ties a finite list into a circular one, or equivalently,
 the infinite repetition of the original list.  It is the identity
 on infinite lists.\par
\subsubsection*{\textbf{Examples}}
\begin{quote}
{\haddockverb\begin{verbatim}
>>> cycle []
*** Exception: Prelude.cycle: empty list

\end{verbatim}}
\end{quote}
\begin{quote}
{\haddockverb\begin{verbatim}
>>> take 10 (cycle [42])
[42,42,42,42,42,42,42,42,42,42]

\end{verbatim}}
\end{quote}
\begin{quote}
{\haddockverb\begin{verbatim}
>>> take 10 (cycle [2, 5, 7])
[2,5,7,2,5,7,2,5,7,2]

\end{verbatim}}
\end{quote}
\begin{quote}
{\haddockverb\begin{verbatim}
>>> take 1 (cycle (42 : undefined))
[42]

\end{verbatim}}
\end{quote}}
\end{haddockdesc}
\subsection{Unfolding}
\begin{haddockdesc}
\item[\begin{tabular}{@{}l}
unfoldr :: (b -> Maybe (a, b)) -> b -> {\char 91}a{\char 93}
\end{tabular}]
{\haddockbegindoc
The \haddockid{unfoldr} function is a `dual' to \haddockid{foldr}: while \haddockid{foldr}
 reduces a list to a summary value, \haddockid{unfoldr} builds a list from
 a seed value.  The function takes the element and returns \haddockid{Nothing}
 if it is done producing the list or returns \haddockid{Just} \haddocktt{(a,b)}, in which
 case, \haddocktt{a} is a prepended to the list and \haddocktt{b} is used as the next
 element in a recursive call.  For example,\par
\begin{quote}
{\haddockverb\begin{verbatim}
iterate f == unfoldr (\x -> Just (x, f x))\end{verbatim}}
\end{quote}
In some cases, \haddockid{unfoldr} can undo a \haddockid{foldr} operation:\par
\begin{quote}
{\haddockverb\begin{verbatim}
unfoldr f' (foldr f z xs) == xs\end{verbatim}}
\end{quote}
if the following holds:\par
\begin{quote}
{\haddockverb\begin{verbatim}
f' (f x y) = Just (x,y)
f' z       = Nothing\end{verbatim}}
\end{quote}
\subsubsection*{\textbf{Laziness}}
\begin{quote}
{\haddockverb\begin{verbatim}
>>> take 1 (unfoldr (\x -> Just (x, undefined)) 'a')
"a"

\end{verbatim}}
\end{quote}
\subsubsection*{\textbf{Examples}}
\begin{quote}
{\haddockverb\begin{verbatim}
>>> unfoldr (\b -> if b == 0 then Nothing else Just (b, b-1)) 10
[10,9,8,7,6,5,4,3,2,1]

\end{verbatim}}
\end{quote}
\begin{quote}
{\haddockverb\begin{verbatim}
>>> take 10 $ unfoldr (\(x, y) -> Just (x, (y, x + y))) (0, 1)
[0,1,1,2,3,5,8,13,21,54]

\end{verbatim}}
\end{quote}}
\end{haddockdesc}
\section{Sublists}
\subsection{Extracting sublists}
\begin{haddockdesc}
\item[\begin{tabular}{@{}l}
take :: Int -> {\char 91}a{\char 93} -> {\char 91}a{\char 93}
\end{tabular}]
{\haddockbegindoc
\haddockid{take} \haddocktt{n}, applied to a list \haddocktt{xs}, returns the prefix of \haddocktt{xs}
 of length \haddocktt{n}, or \haddocktt{xs} itself if \haddocktt{n >= \haddockid{length} xs}.\par
It is an instance of the more general \haddockid{genericTake},
 in which \haddocktt{n} may be of any integral type.\par
\subsubsection*{\textbf{Laziness}}
\begin{quote}
{\haddockverb\begin{verbatim}
>>> take 0 undefined
[]

>>> take 2 (1 : 2 : undefined)
[1,2]

\end{verbatim}}
\end{quote}
\subsubsection*{\textbf{Examples}}
\begin{quote}
{\haddockverb\begin{verbatim}
>>> take 5 "Hello World!"
"Hello"

\end{verbatim}}
\end{quote}
\begin{quote}
{\haddockverb\begin{verbatim}
>>> take 3 [1,2,3,4,5]
[1,2,3]

\end{verbatim}}
\end{quote}
\begin{quote}
{\haddockverb\begin{verbatim}
>>> take 3 [1,2]
[1,2]

\end{verbatim}}
\end{quote}
\begin{quote}
{\haddockverb\begin{verbatim}
>>> take 3 []
[]

\end{verbatim}}
\end{quote}
\begin{quote}
{\haddockverb\begin{verbatim}
>>> take (-1) [1,2]
[]

\end{verbatim}}
\end{quote}
\begin{quote}
{\haddockverb\begin{verbatim}
>>> take 0 [1,2]
[]

\end{verbatim}}
\end{quote}}
\end{haddockdesc}
\begin{haddockdesc}
\item[\begin{tabular}{@{}l}
drop :: Int -> {\char 91}a{\char 93} -> {\char 91}a{\char 93}
\end{tabular}]
{\haddockbegindoc
\haddockid{drop} \haddocktt{n xs} returns the suffix of \haddocktt{xs}
 after the first \haddocktt{n} elements, or \haddocktt{{\char 91}{\char 93}} if \haddocktt{n >= \haddockid{length} xs}.\par
It is an instance of the more general \haddockid{genericDrop},
 in which \haddocktt{n} may be of any integral type.\par
\subsubsection*{\textbf{Examples}}
\begin{quote}
{\haddockverb\begin{verbatim}
>>> drop 6 "Hello World!"
"World!"

\end{verbatim}}
\end{quote}
\begin{quote}
{\haddockverb\begin{verbatim}
>>> drop 3 [1,2,3,4,5]
[4,5]

\end{verbatim}}
\end{quote}
\begin{quote}
{\haddockverb\begin{verbatim}
>>> drop 3 [1,2]
[]

\end{verbatim}}
\end{quote}
\begin{quote}
{\haddockverb\begin{verbatim}
>>> drop 3 []
[]

\end{verbatim}}
\end{quote}
\begin{quote}
{\haddockverb\begin{verbatim}
>>> drop (-1) [1,2]
[1,2]

\end{verbatim}}
\end{quote}
\begin{quote}
{\haddockverb\begin{verbatim}
>>> drop 0 [1,2]
[1,2]

\end{verbatim}}
\end{quote}}
\end{haddockdesc}
\begin{haddockdesc}
\item[\begin{tabular}{@{}l}
splitAt :: Int -> {\char 91}a{\char 93} -> ({\char 91}a{\char 93}, {\char 91}a{\char 93})
\end{tabular}]
{\haddockbegindoc
\haddockid{splitAt} \haddocktt{n xs} returns a tuple where first element is \haddocktt{xs} prefix of
 length \haddocktt{n} and second element is the remainder of the list:\par
\haddockid{splitAt} is an instance of the more general \haddockid{genericSplitAt},
 in which \haddocktt{n} may be of any integral type.\par
\subsubsection*{\textbf{Laziness}}
It is equivalent to \haddocktt{(\haddockid{take} n xs, \haddockid{drop} n xs)}
 unless \haddocktt{n} is \haddocktt{{\char '137}|{\char '137}}:
 \haddocktt{splitAt {\char '137}|{\char '137} xs = {\char '137}|{\char '137}}, not \haddocktt{({\char '137}|{\char '137}, {\char '137}|{\char '137})}).\par
The first component of the tuple is produced lazily:\par
\begin{quote}
{\haddockverb\begin{verbatim}
>>> fst (splitAt 0 undefined)
[]

\end{verbatim}}
\end{quote}
\begin{quote}
{\haddockverb\begin{verbatim}
>>> take 1 (fst (splitAt 10 (1 : undefined)))
[1]

\end{verbatim}}
\end{quote}
\subsubsection*{\textbf{Examples}}
\begin{quote}
{\haddockverb\begin{verbatim}
>>> splitAt 6 "Hello World!"
("Hello ","World!")

\end{verbatim}}
\end{quote}
\begin{quote}
{\haddockverb\begin{verbatim}
>>> splitAt 3 [1,2,3,4,5]
([1,2,3],[4,5])

\end{verbatim}}
\end{quote}
\begin{quote}
{\haddockverb\begin{verbatim}
>>> splitAt 1 [1,2,3]
([1],[2,3])

\end{verbatim}}
\end{quote}
\begin{quote}
{\haddockverb\begin{verbatim}
>>> splitAt 3 [1,2,3]
([1,2,3],[])

\end{verbatim}}
\end{quote}
\begin{quote}
{\haddockverb\begin{verbatim}
>>> splitAt 4 [1,2,3]
([1,2,3],[])

\end{verbatim}}
\end{quote}
\begin{quote}
{\haddockverb\begin{verbatim}
>>> splitAt 0 [1,2,3]
([],[1,2,3])

\end{verbatim}}
\end{quote}
\begin{quote}
{\haddockverb\begin{verbatim}
>>> splitAt (-1) [1,2,3]
([],[1,2,3])

\end{verbatim}}
\end{quote}}
\end{haddockdesc}
\begin{haddockdesc}
\item[\begin{tabular}{@{}l}
takeWhile :: (a -> Bool) -> {\char 91}a{\char 93} -> {\char 91}a{\char 93}
\end{tabular}]
{\haddockbegindoc
\haddockid{takeWhile}, applied to a predicate \haddocktt{p} and a list \haddocktt{xs}, returns the
 longest prefix (possibly empty) of \haddocktt{xs} of elements that satisfy \haddocktt{p}.\par
\subsubsection*{\textbf{Laziness}}
\begin{quote}
{\haddockverb\begin{verbatim}
>>> takeWhile (const False) undefined
*** Exception: Prelude.undefined

\end{verbatim}}
\end{quote}
\begin{quote}
{\haddockverb\begin{verbatim}
>>> takeWhile (const False) (undefined : undefined)
[]

\end{verbatim}}
\end{quote}
\begin{quote}
{\haddockverb\begin{verbatim}
>>> take 1 (takeWhile (const True) (1 : undefined))
[1]

\end{verbatim}}
\end{quote}
\subsubsection*{\textbf{Examples}}
\begin{quote}
{\haddockverb\begin{verbatim}
>>> takeWhile (< 3) [1,2,3,4,1,2,3,4]
[1,2]

\end{verbatim}}
\end{quote}
\begin{quote}
{\haddockverb\begin{verbatim}
>>> takeWhile (< 9) [1,2,3]
[1,2,3]

\end{verbatim}}
\end{quote}
\begin{quote}
{\haddockverb\begin{verbatim}
>>> takeWhile (< 0) [1,2,3]
[]

\end{verbatim}}
\end{quote}}
\end{haddockdesc}
\begin{haddockdesc}
\item[\begin{tabular}{@{}l}
dropWhile :: (a -> Bool) -> {\char 91}a{\char 93} -> {\char 91}a{\char 93}
\end{tabular}]
{\haddockbegindoc
\haddockid{dropWhile} \haddocktt{p xs} returns the suffix remaining after \haddockid{takeWhile} \haddocktt{p xs}.\par
\subsubsection*{\textbf{Examples}}
\begin{quote}
{\haddockverb\begin{verbatim}
>>> dropWhile (< 3) [1,2,3,4,5,1,2,3]
[3,4,5,1,2,3]

\end{verbatim}}
\end{quote}
\begin{quote}
{\haddockverb\begin{verbatim}
>>> dropWhile (< 9) [1,2,3]
[]

\end{verbatim}}
\end{quote}
\begin{quote}
{\haddockverb\begin{verbatim}
>>> dropWhile (< 0) [1,2,3]
[1,2,3]

\end{verbatim}}
\end{quote}}
\end{haddockdesc}
\begin{haddockdesc}
\item[\begin{tabular}{@{}l}
span :: (a -> Bool) -> {\char 91}a{\char 93} -> ({\char 91}a{\char 93}, {\char 91}a{\char 93})
\end{tabular}]
{\haddockbegindoc
\haddockid{span}, applied to a predicate \haddocktt{p} and a list \haddocktt{xs}, returns a tuple where
 first element is the longest prefix (possibly empty) of \haddocktt{xs} of elements that
 satisfy \haddocktt{p} and second element is the remainder of the list:\par
\haddockid{span} \haddocktt{p xs} is equivalent to \haddocktt{(\haddockid{takeWhile} p xs, \haddockid{dropWhile} p xs)}, even if \haddocktt{p} is \haddocktt{{\char '137}|{\char '137}}.\par
\subsubsection*{\textbf{Laziness}}
\begin{quote}
{\haddockverb\begin{verbatim}
>>> span undefined []
([],[])

>>> fst (span (const False) undefined)
*** Exception: Prelude.undefined

>>> fst (span (const False) (undefined : undefined))
[]

>>> take 1 (fst (span (const True) (1 : undefined)))
[1]

\end{verbatim}}
\end{quote}
\haddockid{span} produces the first component of the tuple lazily:\par
\begin{quote}
{\haddockverb\begin{verbatim}
>>> take 10 (fst (span (const True) [1..]))
[1,2,3,4,5,6,7,8,9,10]

\end{verbatim}}
\end{quote}
\subsubsection*{\textbf{Examples}}
\begin{quote}
{\haddockverb\begin{verbatim}
>>> span (< 3) [1,2,3,4,1,2,3,4]
([1,2],[3,4,1,2,3,4])

\end{verbatim}}
\end{quote}
\begin{quote}
{\haddockverb\begin{verbatim}
>>> span (< 9) [1,2,3]
([1,2,3],[])

\end{verbatim}}
\end{quote}
\begin{quote}
{\haddockverb\begin{verbatim}
>>> span (< 0) [1,2,3]
([],[1,2,3])

\end{verbatim}}
\end{quote}}
\end{haddockdesc}
\begin{haddockdesc}
\item[\begin{tabular}{@{}l}
break :: (a -> Bool) -> {\char 91}a{\char 93} -> ({\char 91}a{\char 93}, {\char 91}a{\char 93})
\end{tabular}]
{\haddockbegindoc
\haddockid{break}, applied to a predicate \haddocktt{p} and a list \haddocktt{xs}, returns a tuple where
 first element is longest prefix (possibly empty) of \haddocktt{xs} of elements that
 \emph{do not satisfy} \haddocktt{p} and second element is the remainder of the list:\par
\haddockid{break} \haddocktt{p} is equivalent to \haddocktt{\haddockid{span} (\haddockid{not} . p)}
 and consequently to \haddocktt{(\haddockid{takeWhile} (\haddockid{not} . p) xs, \haddockid{dropWhile} (\haddockid{not} . p) xs)},
 even if \haddocktt{p} is \haddocktt{{\char '137}|{\char '137}}.\par
\subsubsection*{\textbf{Laziness}}
\begin{quote}
{\haddockverb\begin{verbatim}
>>> break undefined []
([],[])

\end{verbatim}}
\end{quote}
\begin{quote}
{\haddockverb\begin{verbatim}
>>> fst (break (const True) undefined)
*** Exception: Prelude.undefined

\end{verbatim}}
\end{quote}
\begin{quote}
{\haddockverb\begin{verbatim}
>>> fst (break (const True) (undefined : undefined))
[]

\end{verbatim}}
\end{quote}
\begin{quote}
{\haddockverb\begin{verbatim}
>>> take 1 (fst (break (const False) (1 : undefined)))
[1]

\end{verbatim}}
\end{quote}
\haddockid{break} produces the first component of the tuple lazily:\par
\begin{quote}
{\haddockverb\begin{verbatim}
>>> take 10 (fst (break (const False) [1..]))
[1,2,3,4,5,6,7,8,9,10]

\end{verbatim}}
\end{quote}
\subsubsection*{\textbf{Examples}}
\begin{quote}
{\haddockverb\begin{verbatim}
>>> break (> 3) [1,2,3,4,1,2,3,4]
([1,2,3],[4,1,2,3,4])

\end{verbatim}}
\end{quote}
\begin{quote}
{\haddockverb\begin{verbatim}
>>> break (< 9) [1,2,3]
([],[1,2,3])

\end{verbatim}}
\end{quote}
\begin{quote}
{\haddockverb\begin{verbatim}
>>> break (> 9) [1,2,3]
([1,2,3],[])

\end{verbatim}}
\end{quote}}
\end{haddockdesc}
\begin{haddockdesc}
\item[\begin{tabular}{@{}l}
stripPrefix :: Eq a => {\char 91}a{\char 93} -> {\char 91}a{\char 93} -> Maybe {\char 91}a{\char 93}
\end{tabular}]
{\haddockbegindoc
\(\mathcal{O}(\min(m,n))\). The \haddockid{stripPrefix} function drops the given
 prefix from a list. It returns \haddockid{Nothing} if the list did not start with the
 prefix given, or \haddockid{Just} the list after the prefix, if it does.\par
\subsubsection*{\textbf{Examples}}
\begin{quote}
{\haddockverb\begin{verbatim}
>>> stripPrefix "foo" "foobar"
Just "bar"

\end{verbatim}}
\end{quote}
\begin{quote}
{\haddockverb\begin{verbatim}
>>> stripPrefix "foo" "foo"
Just ""

\end{verbatim}}
\end{quote}
\begin{quote}
{\haddockverb\begin{verbatim}
>>> stripPrefix "foo" "barfoo"
Nothing

\end{verbatim}}
\end{quote}
\begin{quote}
{\haddockverb\begin{verbatim}
>>> stripPrefix "foo" "barfoobaz"
Nothing

\end{verbatim}}
\end{quote}}
\end{haddockdesc}
\begin{haddockdesc}
\item[\begin{tabular}{@{}l}
group :: Eq a => {\char 91}a{\char 93} -> {\char 91}{\char 91}a{\char 93}{\char 93}
\end{tabular}]
{\haddockbegindoc
The \haddockid{group} function takes a list and returns a list of lists such
 that the concatenation of the result is equal to the argument.  Moreover,
 each sublist in the result is non-empty, all elements are equal to the
 first one, and consecutive equal elements of the input end up in the
 same element of the output list.\par
\haddockid{group} is a special case of \haddockid{groupBy}, which allows the programmer to supply
 their own equality test.\par
It's often preferable to use \haddocktt{Data.List.NonEmpty.}\haddockid{group},
 which provides type-level guarantees of non-emptiness of inner lists.
 A common idiom to squash repeating elements \haddockid{map} \haddockid{head} \haddockid{.} \haddockid{group}
 is better served by
 \haddockid{map} \haddocktt{Data.List.NonEmpty.}\haddockid{head} \haddockid{.} \haddocktt{Data.List.NonEmpty.}\haddockid{group}
 because it avoids partial functions.\par
\subsubsection*{\textbf{Examples}}
\begin{quote}
{\haddockverb\begin{verbatim}
>>> group "Mississippi"
["M","i","ss","i","ss","i","pp","i"]

\end{verbatim}}
\end{quote}
\begin{quote}
{\haddockverb\begin{verbatim}
>>> group [1, 1, 1, 2, 2, 3, 4, 5, 5]
[[1,1,1],[2,2],[3],[4],[5,5]]

\end{verbatim}}
\end{quote}}
\end{haddockdesc}
\begin{haddockdesc}
\item[\begin{tabular}{@{}l}
inits :: {\char 91}a{\char 93} -> {\char 91}{\char 91}a{\char 93}{\char 93}
\end{tabular}]
{\haddockbegindoc
The \haddockid{inits} function returns all initial segments of the argument,
 shortest first.\par
\haddockid{inits} is semantically equivalent to \haddocktt{\haddockid{map} \haddockid{reverse} . \haddockid{scanl} (\haddockid{flip} (:)) {\char 91}{\char 93}},
 but under the hood uses a queue to amortize costs of \haddockid{reverse}.\par
\subsubsection*{\textbf{Laziness}}
Note that \haddockid{inits} has the following strictness property:
 \haddocktt{inits (xs ++ {\char '137}|{\char '137}) = inits xs ++ {\char '137}|{\char '137}}\par
In particular,
 \haddocktt{inits {\char '137}|{\char '137} = {\char 91}{\char 93} : {\char '137}|{\char '137}}\par
\subsubsection*{\textbf{Examples}}
\begin{quote}
{\haddockverb\begin{verbatim}
>>> inits "abc"
["","a","ab","abc"]

\end{verbatim}}
\end{quote}
\begin{quote}
{\haddockverb\begin{verbatim}
>>> inits []
[[]]

\end{verbatim}}
\end{quote}
inits is productive on infinite lists:\par
\begin{quote}
{\haddockverb\begin{verbatim}
>>> take 5 $ inits [1..]
[[],[1],[1,2],[1,2,3],[1,2,3,4]]

\end{verbatim}}
\end{quote}}
\end{haddockdesc}
\begin{haddockdesc}
\item[\begin{tabular}{@{}l}
tails :: {\char 91}a{\char 93} -> {\char 91}{\char 91}a{\char 93}{\char 93}
\end{tabular}]
{\haddockbegindoc
\(\mathcal{O}(n)\). The \haddockid{tails} function returns all final segments of the
 argument, longest first.\par
\subsubsection*{\textbf{Laziness}}
Note that \haddockid{tails} has the following strictness property:
 \haddocktt{tails {\char '137}|{\char '137} = {\char '137}|{\char '137} : {\char '137}|{\char '137}}\par
\begin{quote}
{\haddockverb\begin{verbatim}
>>> tails undefined
[*** Exception: Prelude.undefined

\end{verbatim}}
\end{quote}
\begin{quote}
{\haddockverb\begin{verbatim}
>>> drop 1 (tails [undefined, 1, 2])
[[1, 2], [2], []]

\end{verbatim}}
\end{quote}
\subsubsection*{\textbf{Examples}}
\begin{quote}
{\haddockverb\begin{verbatim}
>>> tails "abc"
["abc","bc","c",""]

\end{verbatim}}
\end{quote}
\begin{quote}
{\haddockverb\begin{verbatim}
>>> tails [1, 2, 3]
[[1,2,3],[2,3],[3],[]]

\end{verbatim}}
\end{quote}
\begin{quote}
{\haddockverb\begin{verbatim}
>>> tails []
[[]]

\end{verbatim}}
\end{quote}}
\end{haddockdesc}
\subsection{Predicates}
\begin{haddockdesc}
\item[\begin{tabular}{@{}l}
isPrefixOf :: Eq a => {\char 91}a{\char 93} -> {\char 91}a{\char 93} -> Bool
\end{tabular}]
{\haddockbegindoc
\(\mathcal{O}(\min(m,n))\). The \haddockid{isPrefixOf} function takes two lists and
 returns \haddockid{True} iff the first list is a prefix of the second.\par
\subsubsection*{\textbf{Examples}}
\begin{quote}
{\haddockverb\begin{verbatim}
>>> "Hello" `isPrefixOf` "Hello World!"
True

\end{verbatim}}
\end{quote}
\begin{quote}
{\haddockverb\begin{verbatim}
>>> "Hello" `isPrefixOf` "Wello Horld!"
False

\end{verbatim}}
\end{quote}
For the result to be \haddockid{True}, the first list must be finite;
 \haddockid{False}, however, results from any mismatch:\par
\begin{quote}
{\haddockverb\begin{verbatim}
>>> [0..] `isPrefixOf` [1..]
False

\end{verbatim}}
\end{quote}
\begin{quote}
{\haddockverb\begin{verbatim}
>>> [0..] `isPrefixOf` [0..99]
False

\end{verbatim}}
\end{quote}
\begin{quote}
{\haddockverb\begin{verbatim}
>>> [0..99] `isPrefixOf` [0..]
True

\end{verbatim}}
\end{quote}
\begin{quote}
{\haddockverb\begin{verbatim}
>>> [0..] `isPrefixOf` [0..]
* Hangs forever *

\end{verbatim}}
\end{quote}
\haddockid{isPrefixOf} shortcuts when the first argument is empty:\par
\begin{quote}
{\haddockverb\begin{verbatim}
>>> isPrefixOf [] undefined
True

\end{verbatim}}
\end{quote}}
\end{haddockdesc}
\begin{haddockdesc}
\item[\begin{tabular}{@{}l}
isSuffixOf :: Eq a => {\char 91}a{\char 93} -> {\char 91}a{\char 93} -> Bool
\end{tabular}]
{\haddockbegindoc
The \haddockid{isSuffixOf} function takes two lists and returns \haddockid{True} iff
 the first list is a suffix of the second.\par
\subsubsection*{\textbf{Examples}}
\begin{quote}
{\haddockverb\begin{verbatim}
>>> "ld!" `isSuffixOf` "Hello World!"
True

\end{verbatim}}
\end{quote}
\begin{quote}
{\haddockverb\begin{verbatim}
>>> "World" `isSuffixOf` "Hello World!"
False

\end{verbatim}}
\end{quote}
The second list must be finite; however the first list may be infinite:\par
\begin{quote}
{\haddockverb\begin{verbatim}
>>> [0..] `isSuffixOf` [0..99]
False

\end{verbatim}}
\end{quote}
\begin{quote}
{\haddockverb\begin{verbatim}
>>> [0..99] `isSuffixOf` [0..]
* Hangs forever *

\end{verbatim}}
\end{quote}}
\end{haddockdesc}
\begin{haddockdesc}
\item[\begin{tabular}{@{}l}
isInfixOf :: Eq a => {\char 91}a{\char 93} -> {\char 91}a{\char 93} -> Bool
\end{tabular}]
{\haddockbegindoc
The \haddockid{isInfixOf} function takes two lists and returns \haddockid{True}
 iff the first list is contained, wholly and intact,
 anywhere within the second.\par
\subsubsection*{\textbf{Examples}}
\begin{quote}
{\haddockverb\begin{verbatim}
>>> isInfixOf "Haskell" "I really like Haskell."
True

\end{verbatim}}
\end{quote}
\begin{quote}
{\haddockverb\begin{verbatim}
>>> isInfixOf "Ial" "I really like Haskell."
False

\end{verbatim}}
\end{quote}
For the result to be \haddockid{True}, the first list must be finite;
 for the result to be \haddockid{False}, the second list must be finite:\par
\begin{quote}
{\haddockverb\begin{verbatim}
>>> [20..50] `isInfixOf` [0..]
True

\end{verbatim}}
\end{quote}
\begin{quote}
{\haddockverb\begin{verbatim}
>>> [0..] `isInfixOf` [20..50]
False

\end{verbatim}}
\end{quote}
\begin{quote}
{\haddockverb\begin{verbatim}
>>> [0..] `isInfixOf` [0..]
* Hangs forever *

\end{verbatim}}
\end{quote}}
\end{haddockdesc}
\section{Searching lists}
\subsection{Searching by equality}
\begin{haddockdesc}
\item[\begin{tabular}{@{}l}
elem :: (Foldable t, Eq a) => a -> t a -> Bool
\end{tabular}]
{\haddockbegindoc
Does the element occur in the structure?\par
Note: \haddockid{elem} is often used in infix form.\par
\subsubsection*{\textbf{Examples}}
Basic usage:\par
\begin{quote}
{\haddockverb\begin{verbatim}
>>> 3 `elem` []
False

\end{verbatim}}
\end{quote}
\begin{quote}
{\haddockverb\begin{verbatim}
>>> 3 `elem` [1,2]
False

\end{verbatim}}
\end{quote}
\begin{quote}
{\haddockverb\begin{verbatim}
>>> 3 `elem` [1,2,3,4,5]
True

\end{verbatim}}
\end{quote}
For infinite structures, the default implementation of \haddockid{elem}
 terminates if the sought-after value exists at a finite distance
 from the left side of the structure:\par
\begin{quote}
{\haddockverb\begin{verbatim}
>>> 3 `elem` [1..]
True

\end{verbatim}}
\end{quote}
\begin{quote}
{\haddockverb\begin{verbatim}
>>> 3 `elem` ([4..] ++ [3])
* Hangs forever *

\end{verbatim}}
\end{quote}}
\end{haddockdesc}
\begin{haddockdesc}
\item[\begin{tabular}{@{}l}
notElem :: (Foldable t, Eq a) => a -> t a -> Bool
\end{tabular}]
{\haddockbegindoc
\haddockid{notElem} is the negation of \haddockid{elem}.\par
\subsubsection*{\textbf{Examples}}
Basic usage:\par
\begin{quote}
{\haddockverb\begin{verbatim}
>>> 3 `notElem` []
True

\end{verbatim}}
\end{quote}
\begin{quote}
{\haddockverb\begin{verbatim}
>>> 3 `notElem` [1,2]
True

\end{verbatim}}
\end{quote}
\begin{quote}
{\haddockverb\begin{verbatim}
>>> 3 `notElem` [1,2,3,4,5]
False

\end{verbatim}}
\end{quote}
For infinite structures, \haddockid{notElem} terminates if the value exists at a
 finite distance from the left side of the structure:\par
\begin{quote}
{\haddockverb\begin{verbatim}
>>> 3 `notElem` [1..]
False

\end{verbatim}}
\end{quote}
\begin{quote}
{\haddockverb\begin{verbatim}
>>> 3 `notElem` ([4..] ++ [3])
* Hangs forever *

\end{verbatim}}
\end{quote}}
\end{haddockdesc}
\begin{haddockdesc}
\item[\begin{tabular}{@{}l}
lookup :: Eq a => a -> {\char 91}(a, b){\char 93} -> Maybe b
\end{tabular}]
{\haddockbegindoc
\(\mathcal{O}(n)\). \haddockid{lookup} \haddocktt{key assocs} looks up a key in an association
 list.
 For the result to be \haddockid{Nothing}, the list must be finite.\par
\subsubsection*{\textbf{Examples}}
\begin{quote}
{\haddockverb\begin{verbatim}
>>> lookup 2 []
Nothing

\end{verbatim}}
\end{quote}
\begin{quote}
{\haddockverb\begin{verbatim}
>>> lookup 2 [(1, "first")]
Nothing

\end{verbatim}}
\end{quote}
\begin{quote}
{\haddockverb\begin{verbatim}
>>> lookup 2 [(1, "first"), (2, "second"), (3, "third")]
Just "second"

\end{verbatim}}
\end{quote}}
\end{haddockdesc}
\subsection{Searching with a predicate}
\begin{haddockdesc}
\item[\begin{tabular}{@{}l}
find :: Foldable t => (a -> Bool) -> t a -> Maybe a
\end{tabular}]
{\haddockbegindoc
The \haddockid{find} function takes a predicate and a structure and returns
 the leftmost element of the structure matching the predicate, or
 \haddockid{Nothing} if there is no such element.\par
\subsubsection*{\textbf{Examples}}
Basic usage:\par
\begin{quote}
{\haddockverb\begin{verbatim}
>>> find (> 42) [0, 5..]
Just 45

\end{verbatim}}
\end{quote}
\begin{quote}
{\haddockverb\begin{verbatim}
>>> find (> 12) [1..7]
Nothing

\end{verbatim}}
\end{quote}}
\end{haddockdesc}
\begin{haddockdesc}
\item[\begin{tabular}{@{}l}
filter :: (a -> Bool) -> {\char 91}a{\char 93} -> {\char 91}a{\char 93}
\end{tabular}]
{\haddockbegindoc
\(\mathcal{O}(n)\). \haddockid{filter}, applied to a predicate and a list, returns
 the list of those elements that satisfy the predicate; i.e.,\par
\begin{quote}
{\haddockverb\begin{verbatim}
filter p xs = [ x | x <- xs, p x]\end{verbatim}}
\end{quote}
\subsubsection*{\textbf{Examples}}
\begin{quote}
{\haddockverb\begin{verbatim}
>>> filter odd [1, 2, 3]
[1,3]

\end{verbatim}}
\end{quote}
\begin{quote}
{\haddockverb\begin{verbatim}
>>> filter (\l -> length l > 3) ["Hello", ", ", "World", "!"]
["Hello","World"]

\end{verbatim}}
\end{quote}
\begin{quote}
{\haddockverb\begin{verbatim}
>>> filter (/= 3) [1, 2, 3, 4, 3, 2, 1]
[1,2,4,2,1]

\end{verbatim}}
\end{quote}}
\end{haddockdesc}
\begin{haddockdesc}
\item[\begin{tabular}{@{}l}
partition :: (a -> Bool) -> {\char 91}a{\char 93} -> ({\char 91}a{\char 93}, {\char 91}a{\char 93})
\end{tabular}]
{\haddockbegindoc
The \haddockid{partition} function takes a predicate and a list, and returns
 the pair of lists of elements which do and do not satisfy the
 predicate, respectively; i.e.,\par
\begin{quote}
{\haddockverb\begin{verbatim}
partition p xs == (filter p xs, filter (not . p) xs)\end{verbatim}}
\end{quote}
\subsubsection*{\textbf{Examples}}
\begin{quote}
{\haddockverb\begin{verbatim}
>>> partition (`elem` "aeiou") "Hello World!"
("eoo","Hll Wrld!")

\end{verbatim}}
\end{quote}
\begin{quote}
{\haddockverb\begin{verbatim}
>>> partition even [1..10]
([2,4,6,8,10],[1,3,5,7,9])

\end{verbatim}}
\end{quote}
\begin{quote}
{\haddockverb\begin{verbatim}
>>> partition (< 5) [1..10]
([1,2,3,4],[5,6,7,8,9,10])

\end{verbatim}}
\end{quote}}
\end{haddockdesc}
\section{Indexing lists}
These functions treat a list \haddocktt{xs} as a indexed collection, with indices
 ranging from \haddocktt{0} to \haddocktt{length xs - 1}.\par
\begin{haddockdesc}
\item[\begin{tabular}{@{}l}
(!!) :: HasCallStack => {\char 91}a{\char 93} -> Int -> a
\end{tabular}]
{\haddockbegindoc
List index (subscript) operator, starting from 0.
 It is an instance of the more general \haddockid{genericIndex},
 which takes an index of any integral type.\par
WARNING: This function is partial, and should only be used if you are
 sure that the indexing will not fail. Otherwise, use \haddockid{!?}.\par
WARNING: This function takes linear time in the index.\par
\subsubsection*{\textbf{Examples}}
\begin{quote}
{\haddockverb\begin{verbatim}
>>> ['a', 'b', 'c'] !! 0
'a'

\end{verbatim}}
\end{quote}
\begin{quote}
{\haddockverb\begin{verbatim}
>>> ['a', 'b', 'c'] !! 2
'c'

\end{verbatim}}
\end{quote}
\begin{quote}
{\haddockverb\begin{verbatim}
>>> ['a', 'b', 'c'] !! 3
*** Exception: Prelude.!!: index too large

\end{verbatim}}
\end{quote}
\begin{quote}
{\haddockverb\begin{verbatim}
>>> ['a', 'b', 'c'] !! (-1)
*** Exception: Prelude.!!: negative index

\end{verbatim}}
\end{quote}}
\end{haddockdesc}
\begin{haddockdesc}
\item[\begin{tabular}{@{}l}
elemIndex :: Eq a => a -> {\char 91}a{\char 93} -> Maybe Int
\end{tabular}]
{\haddockbegindoc
The \haddockid{elemIndex} function returns the index of the first element
 in the given list which is equal (by \haddockid{==}) to the query element,
 or \haddockid{Nothing} if there is no such element.
 For the result to be \haddockid{Nothing}, the list must be finite.\par
\subsubsection*{\textbf{Examples}}
\begin{quote}
{\haddockverb\begin{verbatim}
>>> elemIndex 4 [0..]
Just 4

\end{verbatim}}
\end{quote}
\begin{quote}
{\haddockverb\begin{verbatim}
>>> elemIndex 'o' "haskell"
Nothing

\end{verbatim}}
\end{quote}
\begin{quote}
{\haddockverb\begin{verbatim}
>>> elemIndex 0 [1..]
* hangs forever *

\end{verbatim}}
\end{quote}}
\end{haddockdesc}
\begin{haddockdesc}
\item[\begin{tabular}{@{}l}
elemIndices :: Eq a => a -> {\char 91}a{\char 93} -> {\char 91}Int{\char 93}
\end{tabular}]
{\haddockbegindoc
The \haddockid{elemIndices} function extends \haddockid{elemIndex}, by returning the
 indices of all elements equal to the query element, in ascending order.\par
\subsubsection*{\textbf{Examples}}
\begin{quote}
{\haddockverb\begin{verbatim}
>>> elemIndices 'o' "Hello World"
[4,7]

\end{verbatim}}
\end{quote}
\begin{quote}
{\haddockverb\begin{verbatim}
>>> elemIndices 1 [1, 2, 3, 1, 2, 3]
[0,3]

\end{verbatim}}
\end{quote}}
\end{haddockdesc}
\begin{haddockdesc}
\item[\begin{tabular}{@{}l}
findIndex :: (a -> Bool) -> {\char 91}a{\char 93} -> Maybe Int
\end{tabular}]
{\haddockbegindoc
The \haddockid{findIndex} function takes a predicate and a list and returns
 the index of the first element in the list satisfying the predicate,
 or \haddockid{Nothing} if there is no such element.
 For the result to be \haddockid{Nothing}, the list must be finite.\par
\subsubsection*{\textbf{Examples}}
\begin{quote}
{\haddockverb\begin{verbatim}
>>> findIndex isSpace "Hello World!"
Just 5

\end{verbatim}}
\end{quote}
\begin{quote}
{\haddockverb\begin{verbatim}
>>> findIndex odd [0, 2, 4, 6]
Nothing

\end{verbatim}}
\end{quote}
\begin{quote}
{\haddockverb\begin{verbatim}
>>> findIndex even [1..]
Just 1

\end{verbatim}}
\end{quote}
\begin{quote}
{\haddockverb\begin{verbatim}
>>> findIndex odd [0, 2 ..]
* hangs forever *

\end{verbatim}}
\end{quote}}
\end{haddockdesc}
\begin{haddockdesc}
\item[\begin{tabular}{@{}l}
findIndices :: (a -> Bool) -> {\char 91}a{\char 93} -> {\char 91}Int{\char 93}
\end{tabular}]
{\haddockbegindoc
The \haddockid{findIndices} function extends \haddockid{findIndex}, by returning the
 indices of all elements satisfying the predicate, in ascending order.\par
\subsubsection*{\textbf{Examples}}
\begin{quote}
{\haddockverb\begin{verbatim}
>>> findIndices (`elem` "aeiou") "Hello World!"
[1,4,7]

\end{verbatim}}
\end{quote}
\begin{quote}
{\haddockverb\begin{verbatim}
>>> findIndices (\l -> length l > 3) ["a", "bcde", "fgh", "ijklmnop"]
[1,3]

\end{verbatim}}
\end{quote}}
\end{haddockdesc}
\section{Zipping and unzipping lists}
\begin{haddockdesc}
\item[\begin{tabular}{@{}l}
zip :: {\char 91}a{\char 93} -> {\char 91}b{\char 93} -> {\char 91}(a, b){\char 93}
\end{tabular}]
{\haddockbegindoc
\(\mathcal{O}(\min(m,n))\). \haddockid{zip} takes two lists and returns a list of
 corresponding pairs.\par
\haddockid{zip} is right-lazy:\par
\begin{quote}
{\haddockverb\begin{verbatim}
>>> zip [] undefined
[]

>>> zip undefined []
*** Exception: Prelude.undefined
...

\end{verbatim}}
\end{quote}
\haddockid{zip} is capable of list fusion, but it is restricted to its
 first list argument and its resulting list.\par
\subsubsection*{\textbf{Examples}}
\begin{quote}
{\haddockverb\begin{verbatim}
>>> zip [1, 2, 3] ['a', 'b', 'c']
[(1,'a'),(2,'b'),(3,'c')]

\end{verbatim}}
\end{quote}
If one input list is shorter than the other, excess elements of the longer
 list are discarded, even if one of the lists is infinite:\par
\begin{quote}
{\haddockverb\begin{verbatim}
>>> zip [1] ['a', 'b']
[(1,'a')]

\end{verbatim}}
\end{quote}
\begin{quote}
{\haddockverb\begin{verbatim}
>>> zip [1, 2] ['a']
[(1,'a')]

\end{verbatim}}
\end{quote}
\begin{quote}
{\haddockverb\begin{verbatim}
>>> zip [] [1..]
[]

\end{verbatim}}
\end{quote}
\begin{quote}
{\haddockverb\begin{verbatim}
>>> zip [1..] []
[]

\end{verbatim}}
\end{quote}}
\end{haddockdesc}
\begin{haddockdesc}
\item[\begin{tabular}{@{}l}
zip3 :: {\char 91}a{\char 93} -> {\char 91}b{\char 93} -> {\char 91}c{\char 93} -> {\char 91}(a, b, c){\char 93}
\end{tabular}]
{\haddockbegindoc
\haddockid{zip3} takes three lists and returns a list of triples, analogous to
 \haddockid{zip}.
 It is capable of list fusion, but it is restricted to its
 first list argument and its resulting list.\par}
\end{haddockdesc}
\begin{haddockdesc}
\item[\begin{tabular}{@{}l}
zip4 :: {\char 91}a{\char 93} -> {\char 91}b{\char 93} -> {\char 91}c{\char 93} -> {\char 91}d{\char 93} -> {\char 91}(a, b, c, d){\char 93}
\end{tabular}]
{\haddockbegindoc
The \haddockid{zip4} function takes four lists and returns a list of
 quadruples, analogous to \haddockid{zip}.
 It is capable of list fusion, but it is restricted to its
 first list argument and its resulting list.\par}
\end{haddockdesc}
\begin{haddockdesc}
\item[\begin{tabular}{@{}l}
zip5 :: {\char 91}a{\char 93} -> {\char 91}b{\char 93} -> {\char 91}c{\char 93} -> {\char 91}d{\char 93} -> {\char 91}e{\char 93} -> {\char 91}(a, b, c, d, e){\char 93}
\end{tabular}]
{\haddockbegindoc
The \haddockid{zip5} function takes five lists and returns a list of
 five-tuples, analogous to \haddockid{zip}.
 It is capable of list fusion, but it is restricted to its
 first list argument and its resulting list.\par}
\end{haddockdesc}
\begin{haddockdesc}
\item[\begin{tabular}{@{}l}
zip6 :: {\char 91}a{\char 93} -> {\char 91}b{\char 93} -> {\char 91}c{\char 93} -> {\char 91}d{\char 93} -> {\char 91}e{\char 93} -> {\char 91}f{\char 93} -> {\char 91}(a, b, c, d, e, f){\char 93}
\end{tabular}]
{\haddockbegindoc
The \haddockid{zip6} function takes six lists and returns a list of six-tuples,
 analogous to \haddockid{zip}.
 It is capable of list fusion, but it is restricted to its
 first list argument and its resulting list.\par}
\end{haddockdesc}
\begin{haddockdesc}
\item[\begin{tabular}{@{}l}
zip7 :: {\char 91}a{\char 93} -> {\char 91}b{\char 93} -> {\char 91}c{\char 93} -> {\char 91}d{\char 93} -> {\char 91}e{\char 93} -> {\char 91}f{\char 93} -> {\char 91}g{\char 93} -> {\char 91}(a, b, c, d, e, f, g){\char 93}
\end{tabular}]
{\haddockbegindoc
The \haddockid{zip7} function takes seven lists and returns a list of
 seven-tuples, analogous to \haddockid{zip}.
 It is capable of list fusion, but it is restricted to its
 first list argument and its resulting list.\par}
\end{haddockdesc}
\begin{haddockdesc}
\item[\begin{tabular}{@{}l}
zipWith :: (a -> b -> c) -> {\char 91}a{\char 93} -> {\char 91}b{\char 93} -> {\char 91}c{\char 93}
\end{tabular}]
{\haddockbegindoc
\(\mathcal{O}(\min(m,n))\). \haddockid{zipWith} generalises \haddockid{zip} by zipping with the
 function given as the first argument, instead of a tupling function.\par
\begin{quote}
{\haddockverb\begin{verbatim}
zipWith (,) xs ys == zip xs ys
zipWith f [x1,x2,x3..] [y1,y2,y3..] == [f x1 y1, f x2 y2, f x3 y3..]\end{verbatim}}
\end{quote}
\haddockid{zipWith} is right-lazy:\par
\begin{quote}
{\haddockverb\begin{verbatim}
>>> let f = undefined

>>> zipWith f [] undefined
[]

\end{verbatim}}
\end{quote}
\haddockid{zipWith} is capable of list fusion, but it is restricted to its
 first list argument and its resulting list.\par
\subsubsection*{\textbf{Examples}}
\haddocktt{\haddockid{zipWith} \haddockid{(+)}} can be applied to two lists to produce the list of
 corresponding sums:\par
\begin{quote}
{\haddockverb\begin{verbatim}
>>> zipWith (+) [1, 2, 3] [4, 5, 6]
[5,7,9]

\end{verbatim}}
\end{quote}
\begin{quote}
{\haddockverb\begin{verbatim}
>>> zipWith (++) ["hello ", "foo"] ["world!", "bar"]
["hello world!","foobar"]

\end{verbatim}}
\end{quote}}
\end{haddockdesc}
\begin{haddockdesc}
\item[\begin{tabular}{@{}l}
zipWith3 :: (a -> b -> c -> d) -> {\char 91}a{\char 93} -> {\char 91}b{\char 93} -> {\char 91}c{\char 93} -> {\char 91}d{\char 93}
\end{tabular}]
{\haddockbegindoc
\(\mathcal{O}(\min(l,m,n))\). The \haddockid{zipWith3} function takes a function which combines three
 elements, as well as three lists and returns a list of the function applied
 to corresponding elements, analogous to \haddockid{zipWith}.
 It is capable of list fusion, but it is restricted to its
 first list argument and its resulting list.\par
\begin{quote}
{\haddockverb\begin{verbatim}
zipWith3 (,,) xs ys zs == zip3 xs ys zs
zipWith3 f [x1,x2,x3..] [y1,y2,y3..] [z1,z2,z3..] == [f x1 y1 z1, f x2 y2 z2, f x3 y3 z3..]\end{verbatim}}
\end{quote}
\subsubsection*{\textbf{Examples}}
\begin{quote}
{\haddockverb\begin{verbatim}
>>> zipWith3 (\x y z -> [x, y, z]) "123" "abc" "xyz"
["1ax","2by","3cz"]

\end{verbatim}}
\end{quote}
\begin{quote}
{\haddockverb\begin{verbatim}
>>> zipWith3 (\x y z -> (x * y) + z) [1, 2, 3] [4, 5, 6] [7, 8, 9]
[11,18,27]

\end{verbatim}}
\end{quote}}
\end{haddockdesc}
\begin{haddockdesc}
\item[\begin{tabular}{@{}l}
zipWith4 :: (a -> b -> c -> d -> e) -> {\char 91}a{\char 93} -> {\char 91}b{\char 93} -> {\char 91}c{\char 93} -> {\char 91}d{\char 93} -> {\char 91}e{\char 93}
\end{tabular}]
{\haddockbegindoc
The \haddockid{zipWith4} function takes a function which combines four
 elements, as well as four lists and returns a list of their point-wise
 combination, analogous to \haddockid{zipWith}.
 It is capable of list fusion, but it is restricted to its
 first list argument and its resulting list.\par}
\end{haddockdesc}
\begin{haddockdesc}
\item[\begin{tabular}{@{}l}
zipWith5 :: (a -> b -> c -> d -> e -> f) -> {\char 91}a{\char 93} -> {\char 91}b{\char 93} -> {\char 91}c{\char 93} -> {\char 91}d{\char 93} -> {\char 91}e{\char 93} -> {\char 91}f{\char 93}
\end{tabular}]
{\haddockbegindoc
The \haddockid{zipWith5} function takes a function which combines five
 elements, as well as five lists and returns a list of their point-wise
 combination, analogous to \haddockid{zipWith}.
 It is capable of list fusion, but it is restricted to its
 first list argument and its resulting list.\par}
\end{haddockdesc}
\begin{haddockdesc}
\item[\begin{tabular}{@{}l}
zipWith6 :: (a -> b -> c -> d -> e -> f -> g)\\\ \ \ \ \ \ \ \ \ \ \ \ -> {\char 91}a{\char 93} -> {\char 91}b{\char 93} -> {\char 91}c{\char 93} -> {\char 91}d{\char 93} -> {\char 91}e{\char 93} -> {\char 91}f{\char 93} -> {\char 91}g{\char 93}
\end{tabular}]
{\haddockbegindoc
The \haddockid{zipWith6} function takes a function which combines six
 elements, as well as six lists and returns a list of their point-wise
 combination, analogous to \haddockid{zipWith}.
 It is capable of list fusion, but it is restricted to its
 first list argument and its resulting list.\par}
\end{haddockdesc}
\begin{haddockdesc}
\item[\begin{tabular}{@{}l}
zipWith7 :: (a -> b -> c -> d -> e -> f -> g -> h)\\\ \ \ \ \ \ \ \ \ \ \ \ -> {\char 91}a{\char 93} -> {\char 91}b{\char 93} -> {\char 91}c{\char 93} -> {\char 91}d{\char 93} -> {\char 91}e{\char 93} -> {\char 91}f{\char 93} -> {\char 91}g{\char 93} -> {\char 91}h{\char 93}
\end{tabular}]
{\haddockbegindoc
The \haddockid{zipWith7} function takes a function which combines seven
 elements, as well as seven lists and returns a list of their point-wise
 combination, analogous to \haddockid{zipWith}.
 It is capable of list fusion, but it is restricted to its
 first list argument and its resulting list.\par}
\end{haddockdesc}
\begin{haddockdesc}
\item[\begin{tabular}{@{}l}
unzip :: {\char 91}(a, b){\char 93} -> ({\char 91}a{\char 93}, {\char 91}b{\char 93})
\end{tabular}]
{\haddockbegindoc
\haddockid{unzip} transforms a list of pairs into a list of first components
 and a list of second components.\par
\subsubsection*{\textbf{Examples}}
\begin{quote}
{\haddockverb\begin{verbatim}
>>> unzip []
([],[])

\end{verbatim}}
\end{quote}
\begin{quote}
{\haddockverb\begin{verbatim}
>>> unzip [(1, 'a'), (2, 'b')]
([1,2],"ab")

\end{verbatim}}
\end{quote}}
\end{haddockdesc}
\begin{haddockdesc}
\item[\begin{tabular}{@{}l}
unzip3 :: {\char 91}(a, b, c){\char 93} -> ({\char 91}a{\char 93}, {\char 91}b{\char 93}, {\char 91}c{\char 93})
\end{tabular}]
{\haddockbegindoc
The \haddockid{unzip3} function takes a list of triples and returns three
 lists of the respective components, analogous to \haddockid{unzip}.\par
\subsubsection*{\textbf{Examples}}
\begin{quote}
{\haddockverb\begin{verbatim}
>>> unzip3 []
([],[],[])

\end{verbatim}}
\end{quote}
\begin{quote}
{\haddockverb\begin{verbatim}
>>> unzip3 [(1, 'a', True), (2, 'b', False)]
([1,2],"ab",[True,False])

\end{verbatim}}
\end{quote}}
\end{haddockdesc}
\begin{haddockdesc}
\item[\begin{tabular}{@{}l}
unzip4 :: {\char 91}(a, b, c, d){\char 93} -> ({\char 91}a{\char 93}, {\char 91}b{\char 93}, {\char 91}c{\char 93}, {\char 91}d{\char 93})
\end{tabular}]
{\haddockbegindoc
The \haddockid{unzip4} function takes a list of quadruples and returns four
 lists, analogous to \haddockid{unzip}.\par}
\end{haddockdesc}
\begin{haddockdesc}
\item[\begin{tabular}{@{}l}
unzip5 :: {\char 91}(a, b, c, d, e){\char 93} -> ({\char 91}a{\char 93}, {\char 91}b{\char 93}, {\char 91}c{\char 93}, {\char 91}d{\char 93}, {\char 91}e{\char 93})
\end{tabular}]
{\haddockbegindoc
The \haddockid{unzip5} function takes a list of five-tuples and returns five
 lists, analogous to \haddockid{unzip}.\par}
\end{haddockdesc}
\begin{haddockdesc}
\item[\begin{tabular}{@{}l}
unzip6 :: {\char 91}(a, b, c, d, e, f){\char 93} -> ({\char 91}a{\char 93}, {\char 91}b{\char 93}, {\char 91}c{\char 93}, {\char 91}d{\char 93}, {\char 91}e{\char 93}, {\char 91}f{\char 93})
\end{tabular}]
{\haddockbegindoc
The \haddockid{unzip6} function takes a list of six-tuples and returns six
 lists, analogous to \haddockid{unzip}.\par}
\end{haddockdesc}
\begin{haddockdesc}
\item[\begin{tabular}{@{}l}
unzip7 :: {\char 91}(a, b, c, d, e, f, g){\char 93} -> ({\char 91}a{\char 93}, {\char 91}b{\char 93}, {\char 91}c{\char 93}, {\char 91}d{\char 93}, {\char 91}e{\char 93}, {\char 91}f{\char 93}, {\char 91}g{\char 93})
\end{tabular}]
{\haddockbegindoc
The \haddockid{unzip7} function takes a list of seven-tuples and returns
 seven lists, analogous to \haddockid{unzip}.\par}
\end{haddockdesc}
\section{Special lists}
\subsection{Functions on strings}
\begin{haddockdesc}
\item[\begin{tabular}{@{}l}
lines :: String -> {\char 91}String{\char 93}
\end{tabular}]
{\haddockbegindoc
Splits the argument into a list of \emph{lines} stripped of their terminating
 \haddocktt{{\char '134}n} characters.  The \haddocktt{{\char '134}n} terminator is optional in a final non-empty
 line of the argument string.\par
When the argument string is empty, or ends in a \haddocktt{{\char '134}n} character, it can be
 recovered by passing the result of \haddockid{lines} to the \haddockid{unlines} function.
 Otherwise, \haddockid{unlines} appends the missing terminating \haddocktt{{\char '134}n}.  This makes
 \haddocktt{unlines . lines} \emph{idempotent}:\par
\begin{quote}
{\haddockverb\begin{verbatim}
(unlines . lines) . (unlines . lines) = (unlines . lines)\end{verbatim}}
\end{quote}
\subsubsection*{\textbf{Examples}}
\begin{quote}
{\haddockverb\begin{verbatim}
>>> lines ""           -- empty input contains no lines
[]

\end{verbatim}}
\end{quote}
\begin{quote}
{\haddockverb\begin{verbatim}
>>> lines "\n"         -- single empty line
[""]

\end{verbatim}}
\end{quote}
\begin{quote}
{\haddockverb\begin{verbatim}
>>> lines "one"        -- single unterminated line
["one"]

\end{verbatim}}
\end{quote}
\begin{quote}
{\haddockverb\begin{verbatim}
>>> lines "one\n"      -- single non-empty line
["one"]

\end{verbatim}}
\end{quote}
\begin{quote}
{\haddockverb\begin{verbatim}
>>> lines "one\n\n"    -- second line is empty
["one",""]

\end{verbatim}}
\end{quote}
\begin{quote}
{\haddockverb\begin{verbatim}
>>> lines "one\ntwo"   -- second line is unterminated
["one","two"]

\end{verbatim}}
\end{quote}
\begin{quote}
{\haddockverb\begin{verbatim}
>>> lines "one\ntwo\n" -- two non-empty lines
["one","two"]

\end{verbatim}}
\end{quote}}
\end{haddockdesc}
\begin{haddockdesc}
\item[\begin{tabular}{@{}l}
words :: String -> {\char 91}String{\char 93}
\end{tabular}]
{\haddockbegindoc
\haddockid{words} breaks a string up into a list of words, which were delimited
 by white space (as defined by \haddockid{isSpace}). This function trims any white spaces
 at the beginning and at the end.\par
\subsubsection*{\textbf{Examples}}
\begin{quote}
{\haddockverb\begin{verbatim}
>>> words "Lorem ipsum\ndolor"
["Lorem","ipsum","dolor"]

\end{verbatim}}
\end{quote}
\begin{quote}
{\haddockverb\begin{verbatim}
>>> words " foo bar "
["foo","bar"]

\end{verbatim}}
\end{quote}}
\end{haddockdesc}
\begin{haddockdesc}
\item[\begin{tabular}{@{}l}
unlines :: {\char 91}String{\char 93} -> String
\end{tabular}]
{\haddockbegindoc
Appends a \haddocktt{{\char '134}n} character to each input string, then concatenates the
 results. Equivalent to \haddocktt{\haddocktt{foldMap} (s -> s \haddockid{++} "{\char '134}n")}.\par
\subsubsection*{\textbf{Examples}}
\begin{quote}
{\haddockverb\begin{verbatim}
>>> unlines ["Hello", "World", "!"]
"Hello\nWorld\n!\n"

\end{verbatim}}
\end{quote}
Note that \haddocktt{\haddockid{unlines} \haddockid{.} \haddockid{lines} \haddockid{/=} \haddockid{id}} when the input is not \haddocktt{{\char '134}n}-terminated:\par
\begin{quote}
{\haddockverb\begin{verbatim}
>>> unlines . lines $ "foo\nbar"
"foo\nbar\n"

\end{verbatim}}
\end{quote}}
\end{haddockdesc}
\begin{haddockdesc}
\item[\begin{tabular}{@{}l}
unwords :: {\char 91}String{\char 93} -> String
\end{tabular}]
{\haddockbegindoc
\haddockid{unwords} joins words with separating spaces (U+0020 SPACE).\par
\haddockid{unwords} is neither left nor right inverse of \haddockid{words}:\par
\begin{quote}
{\haddockverb\begin{verbatim}
>>> words (unwords [" "])
[]

>>> unwords (words "foo\nbar")
"foo bar"

\end{verbatim}}
\end{quote}
\subsubsection*{\textbf{Examples}}
\begin{quote}
{\haddockverb\begin{verbatim}
>>> unwords ["Lorem", "ipsum", "dolor"]
"Lorem ipsum dolor"

\end{verbatim}}
\end{quote}
\begin{quote}
{\haddockverb\begin{verbatim}
>>> unwords ["foo", "bar", "", "baz"]
"foo bar  baz"

\end{verbatim}}
\end{quote}}
\end{haddockdesc}
\subsection{\haddocktt{Set} operations}
\begin{haddockdesc}
\item[\begin{tabular}{@{}l}
nub :: Eq a => {\char 91}a{\char 93} -> {\char 91}a{\char 93}
\end{tabular}]
{\haddockbegindoc
\(\mathcal{O}(n^2)\). The \haddockid{nub} function removes duplicate elements from a
 list. In particular, it keeps only the first occurrence of each element. (The
 name \haddockid{nub} means `essence'.) It is a special case of \haddockid{nubBy}, which allows
 the programmer to supply their own equality test.\par
If there exists \haddocktt{instance Ord a}, it's faster to use \haddocktt{nubOrd} from the \haddocktt{containers} package
 (\href{https://hackage.haskell.org/package/containers/docs/Data-Containers-ListUtils.html#v:nubOrd}{link to the latest online documentation}),
 which takes only \(\mathcal{O}(n \log d)\) time where \haddocktt{d} is the number of
 distinct elements in the list.\par
Another approach to speed up \haddockid{nub} is to use
 \haddockid{map} \haddocktt{Data.List.NonEmpty.}\haddockid{head} . \haddocktt{Data.List.NonEmpty.}\haddockid{group} . \haddockid{sort},
 which takes \(\mathcal{O}(n \log n)\) time, requires \haddocktt{instance Ord a} and doesn't
 preserve the order.\par
\subsubsection*{\textbf{Examples}}
\begin{quote}
{\haddockverb\begin{verbatim}
>>> nub [1,2,3,4,3,2,1,2,4,3,5]
[1,2,3,4,5]

\end{verbatim}}
\end{quote}
\begin{quote}
{\haddockverb\begin{verbatim}
>>> nub "hello, world!"
"helo, wrd!"

\end{verbatim}}
\end{quote}}
\end{haddockdesc}
\begin{haddockdesc}
\item[\begin{tabular}{@{}l}
delete :: Eq a => a -> {\char 91}a{\char 93} -> {\char 91}a{\char 93}
\end{tabular}]
{\haddockbegindoc
\(\mathcal{O}(n)\). \haddockid{delete} \haddocktt{x} removes the first occurrence of \haddocktt{x} from
 its list argument.\par
It is a special case of \haddockid{deleteBy}, which allows the programmer to
 supply their own equality test.\par
\subsubsection*{\textbf{Examples}}
\begin{quote}
{\haddockverb\begin{verbatim}
>>> delete 'a' "banana"
"bnana"

\end{verbatim}}
\end{quote}
\begin{quote}
{\haddockverb\begin{verbatim}
>>> delete "not" ["haskell", "is", "not", "awesome"]
["haskell","is","awesome"]

\end{verbatim}}
\end{quote}}
\end{haddockdesc}
\begin{haddockdesc}
\item[\begin{tabular}{@{}l}
({\char '134}{\char '134}) :: Eq a => {\char 91}a{\char 93} -> {\char 91}a{\char 93} -> {\char 91}a{\char 93}
\end{tabular}]
{\haddockbegindoc
The \haddockid{{\char '134}{\char '134}} function is list difference (non-associative).
 In the result of \haddocktt{xs} \haddockid{{\char '134}{\char '134}} \haddocktt{ys}, the first occurrence of each element of
 \haddocktt{ys} in turn (if any) has been removed from \haddocktt{xs}.  Thus
 \haddocktt{(xs ++ ys) {\char '134}{\char '134} xs == ys}.\par
It is a special case of \haddockid{deleteFirstsBy}, which allows the programmer
 to supply their own equality test.\par
\subsubsection*{\textbf{Examples}}
\begin{quote}
{\haddockverb\begin{verbatim}
>>> "Hello World!" \\ "ell W"
"Hoorld!"

\end{verbatim}}
\end{quote}
The second list must be finite, but the first may be infinite.\par
\begin{quote}
{\haddockverb\begin{verbatim}
>>> take 5 ([0..] \\ [2..4])
[0,1,5,6,7]

\end{verbatim}}
\end{quote}
\begin{quote}
{\haddockverb\begin{verbatim}
>>> take 5 ([0..] \\ [2..])
* Hangs forever *

\end{verbatim}}
\end{quote}}
\end{haddockdesc}
\begin{haddockdesc}
\item[\begin{tabular}{@{}l}
union :: Eq a => {\char 91}a{\char 93} -> {\char 91}a{\char 93} -> {\char 91}a{\char 93}
\end{tabular}]
{\haddockbegindoc
The \haddockid{union} function returns the list union of the two lists.
 It is a special case of \haddockid{unionBy}, which allows the programmer to supply
 their own equality test.\par
\subsubsection*{\textbf{Examples}}
\begin{quote}
{\haddockverb\begin{verbatim}
>>> "dog" `union` "cow"
"dogcw"

\end{verbatim}}
\end{quote}
If equal elements are present in both lists, an element from the first list
 will be used. If the second list contains equal elements, only the first one
 will be retained:\par
\begin{quote}
{\haddockverb\begin{verbatim}
>>> import Data.Semigroup(Arg(..))

>>> union [Arg () "dog"] [Arg () "cow"]
[Arg () "dog"]

>>> union [] [Arg () "dog", Arg () "cow"]
[Arg () "dog"]

\end{verbatim}}
\end{quote}
However if the first list contains duplicates, so will
 the result:\par
\begin{quote}
{\haddockverb\begin{verbatim}
>>> "coot" `union` "duck"
"cootduk"

>>> "duck" `union` "coot"
"duckot"

\end{verbatim}}
\end{quote}
\haddockid{union} is productive even if both arguments are infinite.\par
\begin{quote}
{\haddockverb\begin{verbatim}
>>> [0, 2 ..] `union` [1, 3 ..]
[0,2,4,6,8,10,12..

\end{verbatim}}
\end{quote}}
\end{haddockdesc}
\begin{haddockdesc}
\item[\begin{tabular}{@{}l}
intersect :: Eq a => {\char 91}a{\char 93} -> {\char 91}a{\char 93} -> {\char 91}a{\char 93}
\end{tabular}]
{\haddockbegindoc
The \haddockid{intersect} function takes the list intersection of two lists.
 It is a special case of \haddockid{intersectBy}, which allows the programmer to
 supply their own equality test.\par
\subsubsection*{\textbf{Examples}}
\begin{quote}
{\haddockverb\begin{verbatim}
>>> [1,2,3,4] `intersect` [2,4,6,8]
[2,4]

\end{verbatim}}
\end{quote}
If equal elements are present in both lists, an element from the first list
 will be used, and all duplicates from the second list quashed:\par
\begin{quote}
{\haddockverb\begin{verbatim}
>>> import Data.Semigroup

>>> intersect [Arg () "dog"] [Arg () "cow", Arg () "cat"]
[Arg () "dog"]

\end{verbatim}}
\end{quote}
However if the first list contains duplicates, so will the result.\par
\begin{quote}
{\haddockverb\begin{verbatim}
>>> "coot" `intersect` "heron"
"oo"

>>> "heron" `intersect` "coot"
"o"

\end{verbatim}}
\end{quote}
If the second list is infinite, \haddockid{intersect} either hangs
 or returns its first argument in full. Otherwise if the first list
 is infinite, \haddockid{intersect} might be productive:\par
\begin{quote}
{\haddockverb\begin{verbatim}
>>> intersect [100..] [0..]
[100,101,102,103...

>>> intersect [0] [1..]
* Hangs forever *

>>> intersect [1..] [0]
* Hangs forever *

>>> intersect (cycle [1..3]) [2]
[2,2,2,2...

\end{verbatim}}
\end{quote}}
\end{haddockdesc}
\subsection{Ordered lists}
\begin{haddockdesc}
\item[\begin{tabular}{@{}l}
sort :: Ord a => {\char 91}a{\char 93} -> {\char 91}a{\char 93}
\end{tabular}]
{\haddockbegindoc
The \haddockid{sort} function implements a stable sorting algorithm.
 It is a special case of \haddockid{sortBy}, which allows the programmer to supply
 their own comparison function.\par
Elements are arranged from lowest to highest, keeping duplicates in
 the order they appeared in the input.\par
The argument must be finite.\par
\subsubsection*{\textbf{Examples}}
\begin{quote}
{\haddockverb\begin{verbatim}
>>> sort [1,6,4,3,2,5]
[1,2,3,4,5,6]

\end{verbatim}}
\end{quote}
\begin{quote}
{\haddockverb\begin{verbatim}
>>> sort "haskell"
"aehklls"

\end{verbatim}}
\end{quote}
\begin{quote}
{\haddockverb\begin{verbatim}
>>> import Data.Semigroup(Arg(..))

>>> sort [Arg ":)" 0, Arg ":D" 0, Arg ":)" 1, Arg ":3" 0, Arg ":D" 1]
[Arg ":)" 0,Arg ":)" 1,Arg ":3" 0,Arg ":D" 0,Arg ":D" 1]

\end{verbatim}}
\end{quote}}
\end{haddockdesc}
\begin{haddockdesc}
\item[\begin{tabular}{@{}l}
insert :: Ord a => a -> {\char 91}a{\char 93} -> {\char 91}a{\char 93}
\end{tabular}]
{\haddockbegindoc
\(\mathcal{O}(n)\). The \haddockid{insert} function takes an element and a list and
 inserts the element into the list at the first position where it is less than
 or equal to the next element. In particular, if the list is sorted before the
 call, the result will also be sorted. It is a special case of \haddockid{insertBy},
 which allows the programmer to supply their own comparison function.\par
\subsubsection*{\textbf{Examples}}
\begin{quote}
{\haddockverb\begin{verbatim}
>>> insert (-1) [1, 2, 3]
[-1,1,2,3]

\end{verbatim}}
\end{quote}
\begin{quote}
{\haddockverb\begin{verbatim}
>>> insert 'd' "abcefg"
"abcdefg"

\end{verbatim}}
\end{quote}
\begin{quote}
{\haddockverb\begin{verbatim}
>>> insert 4 [1, 2, 3, 5, 6, 7]
[1,2,3,4,5,6,7]

\end{verbatim}}
\end{quote}}
\end{haddockdesc}
\section{Generalized functions}
\subsection{The \haddocktt{By} operations}
By convention, overloaded functions have a non-overloaded counterpart
 whose name is suffixed with ‘By’.
 *** User-supplied equality (replacing an Eq context)
 | The predicate is assumed to define an equivalence.\par
\begin{haddockdesc}
\item[\begin{tabular}{@{}l}
nubBy :: (a -> a -> Bool) -> {\char 91}a{\char 93} -> {\char 91}a{\char 93}
\end{tabular}]
{\haddockbegindoc
The \haddockid{nubBy} function behaves just like \haddockid{nub}, except it uses a
 user-supplied equality predicate instead of the overloaded \haddockid{(==)}
 function.\par
\subsubsection*{\textbf{Examples}}
\begin{quote}
{\haddockverb\begin{verbatim}
>>> nubBy (\x y -> mod x 3 == mod y 3) [1,2,4,5,6]
[1,2,6]

\end{verbatim}}
\end{quote}
\begin{quote}
{\haddockverb\begin{verbatim}
>>> nubBy (/=) [2, 7, 1, 8, 2, 8, 1, 8, 2, 8]
[2,2,2]

\end{verbatim}}
\end{quote}
\begin{quote}
{\haddockverb\begin{verbatim}
>>> nubBy (>) [1, 2, 3, 2, 1, 5, 4, 5, 3, 2]
[1,2,3,5,5]

\end{verbatim}}
\end{quote}}
\end{haddockdesc}
\begin{haddockdesc}
\item[\begin{tabular}{@{}l}
deleteBy :: (a -> a -> Bool) -> a -> {\char 91}a{\char 93} -> {\char 91}a{\char 93}
\end{tabular}]
{\haddockbegindoc
\(\mathcal{O}(n)\). The \haddockid{deleteBy} function behaves like \haddockid{delete}, but
 takes a user-supplied equality predicate.\par
\subsubsection*{\textbf{Examples}}
\begin{quote}
{\haddockverb\begin{verbatim}
>>> deleteBy (<=) 4 [1..10]
[1,2,3,5,6,7,8,9,10]

\end{verbatim}}
\end{quote}
\begin{quote}
{\haddockverb\begin{verbatim}
>>> deleteBy (/=) 5 [5, 5, 4, 3, 5, 2]
[5,5,3,5,2]

\end{verbatim}}
\end{quote}}
\end{haddockdesc}
\begin{haddockdesc}
\item[\begin{tabular}{@{}l}
deleteFirstsBy :: (a -> a -> Bool) -> {\char 91}a{\char 93} -> {\char 91}a{\char 93} -> {\char 91}a{\char 93}
\end{tabular}]
{\haddockbegindoc
The \haddockid{deleteFirstsBy} function takes a predicate and two lists and
 returns the first list with the first occurrence of each element of
 the second list removed. This is the non-overloaded version of \haddockid{({\char '134}{\char '134})}.\par
\begin{quote}
{\haddockverb\begin{verbatim}
(\\) == deleteFirstsBy (==)\end{verbatim}}
\end{quote}
The second list must be finite, but the first may be infinite.\par
\subsubsection*{\textbf{Examples}}
\begin{quote}
{\haddockverb\begin{verbatim}
>>> deleteFirstsBy (>) [1..10] [3, 4, 5]
[4,5,6,7,8,9,10]

\end{verbatim}}
\end{quote}
\begin{quote}
{\haddockverb\begin{verbatim}
>>> deleteFirstsBy (/=) [1..10] [1, 3, 5]
[4,5,6,7,8,9,10]

\end{verbatim}}
\end{quote}}
\end{haddockdesc}
\begin{haddockdesc}
\item[\begin{tabular}{@{}l}
unionBy :: (a -> a -> Bool) -> {\char 91}a{\char 93} -> {\char 91}a{\char 93} -> {\char 91}a{\char 93}
\end{tabular}]
{\haddockbegindoc
The \haddockid{unionBy} function is the non-overloaded version of \haddockid{union}.
 Both arguments may be infinite.\par
\subsubsection*{\textbf{Examples}}
\begin{quote}
{\haddockverb\begin{verbatim}
>>> unionBy (>) [3, 4, 5] [1, 2, 3, 4, 5, 6]
[3,4,5,4,5,6]

\end{verbatim}}
\end{quote}
\begin{quote}
{\haddockverb\begin{verbatim}
>>> import Data.Semigroup (Arg(..))

>>> unionBy (/=) [Arg () "Saul"] [Arg () "Kim"]
[Arg () "Saul", Arg () "Kim"]

\end{verbatim}}
\end{quote}}
\end{haddockdesc}
\begin{haddockdesc}
\item[\begin{tabular}{@{}l}
intersectBy :: (a -> a -> Bool) -> {\char 91}a{\char 93} -> {\char 91}a{\char 93} -> {\char 91}a{\char 93}
\end{tabular}]
{\haddockbegindoc
The \haddockid{intersectBy} function is the non-overloaded version of \haddockid{intersect}.
 It is productive for infinite arguments only if the first one
 is a subset of the second.\par}
\end{haddockdesc}
\begin{haddockdesc}
\item[\begin{tabular}{@{}l}
groupBy :: (a -> a -> Bool) -> {\char 91}a{\char 93} -> {\char 91}{\char 91}a{\char 93}{\char 93}
\end{tabular}]
{\haddockbegindoc
The \haddockid{groupBy} function is the non-overloaded version of \haddockid{group}.\par
When a supplied relation is not transitive, it is important
 to remember that equality is checked against the first element in the group,
 not against the nearest neighbour:\par
\begin{quote}
{\haddockverb\begin{verbatim}
>>> groupBy (\a b -> b - a < 5) [0..19]
[[0,1,2,3,4],[5,6,7,8,9],[10,11,12,13,14],[15,16,17,18,19]]

\end{verbatim}}
\end{quote}
It's often preferable to use \haddocktt{Data.List.NonEmpty.}\haddockid{groupBy},
 which provides type-level guarantees of non-emptiness of inner lists.\par
\subsubsection*{\textbf{Examples}}
\begin{quote}
{\haddockverb\begin{verbatim}
>>> groupBy (/=) [1, 1, 1, 2, 3, 1, 4, 4, 5]
[[1],[1],[1,2,3],[1,4,4,5]]

\end{verbatim}}
\end{quote}
\begin{quote}
{\haddockverb\begin{verbatim}
>>> groupBy (>) [1, 3, 5, 1, 4, 2, 6, 5, 4]
[[1],[3],[5,1,4,2],[6,5,4]]

\end{verbatim}}
\end{quote}
\begin{quote}
{\haddockverb\begin{verbatim}
>>> groupBy (const not) [True, False, True, False, False, False, True]
[[True,False],[True,False,False,False],[True]]

\end{verbatim}}
\end{quote}}
\end{haddockdesc}
\subsubsection{User-supplied comparison (replacing an Ord context)}
The function is assumed to define a total ordering.\par
\begin{haddockdesc}
\item[\begin{tabular}{@{}l}
sortBy :: (a -> a -> Ordering) -> {\char 91}a{\char 93} -> {\char 91}a{\char 93}
\end{tabular}]
{\haddockbegindoc
The \haddockid{sortBy} function is the non-overloaded version of \haddockid{sort}.
 The argument must be finite.\par
The supplied comparison relation is supposed to be reflexive and antisymmetric,
 otherwise, e. g., for \haddocktt{{\char '137} {\char '137} -> GT}, the ordered list simply does not exist.
 The relation is also expected to be transitive: if it is not then \haddockid{sortBy}
 might fail to find an ordered permutation, even if it exists.\par
\subsubsection*{\textbf{Examples}}
\begin{quote}
{\haddockverb\begin{verbatim}
>>> sortBy (\(a,_) (b,_) -> compare a b) [(2, "world"), (4, "!"), (1, "Hello")]
[(1,"Hello"),(2,"world"),(4,"!")]

\end{verbatim}}
\end{quote}}
\end{haddockdesc}
\begin{haddockdesc}
\item[\begin{tabular}{@{}l}
insertBy :: (a -> a -> Ordering) -> a -> {\char 91}a{\char 93} -> {\char 91}a{\char 93}
\end{tabular}]
{\haddockbegindoc
\(\mathcal{O}(n)\). The non-overloaded version of \haddockid{insert}.\par
\subsubsection*{\textbf{Examples}}
\begin{quote}
{\haddockverb\begin{verbatim}
>>> insertBy (\x y -> compare (length x) (length y)) [1, 2] [[1], [1, 2, 3], [1, 2, 3, 4]]
[[1],[1,2],[1,2,3],[1,2,3,4]]

\end{verbatim}}
\end{quote}}
\end{haddockdesc}
\begin{haddockdesc}
\item[\begin{tabular}{@{}l}
maximumBy :: Foldable t => (a -> a -> Ordering) -> t a -> a
\end{tabular}]
{\haddockbegindoc
The largest element of a non-empty structure with respect to the
 given comparison function.\par
\subsubsection*{\textbf{Examples}}
Basic usage:\par
\begin{quote}
{\haddockverb\begin{verbatim}
>>> maximumBy (compare `on` length) ["Hello", "World", "!", "Longest", "bar"]
"Longest"

\end{verbatim}}
\end{quote}
WARNING: This function is partial for possibly-empty structures like lists.\par}
\end{haddockdesc}
\begin{haddockdesc}
\item[\begin{tabular}{@{}l}
minimumBy :: Foldable t => (a -> a -> Ordering) -> t a -> a
\end{tabular}]
{\haddockbegindoc
The least element of a non-empty structure with respect to the
 given comparison function.\par
\subsubsection*{\textbf{Examples}}
Basic usage:\par
\begin{quote}
{\haddockverb\begin{verbatim}
>>> minimumBy (compare `on` length) ["Hello", "World", "!", "Longest", "bar"]
"!"

\end{verbatim}}
\end{quote}
WARNING: This function is partial for possibly-empty structures like lists.\par}
\end{haddockdesc}
\subsection{The "generic" operations}
The prefix ‘generic’ indicates an overloaded function that is a
 generalized version of a \haddocktt{Prelude} function.\par
\begin{haddockdesc}
\item[\begin{tabular}{@{}l}
genericLength :: Num i => {\char 91}a{\char 93} -> i
\end{tabular}]
{\haddockbegindoc
\(\mathcal{O}(n)\). The \haddockid{genericLength} function is an overloaded version
 of \haddockid{length}. In particular, instead of returning an \haddockid{Int}, it returns any
 type which is an instance of \haddockid{Num}. It is, however, less efficient than
 \haddockid{length}.\par
\subsubsection*{\textbf{Examples}}
\begin{quote}
{\haddockverb\begin{verbatim}
>>> genericLength [1, 2, 3] :: Int
3

>>> genericLength [1, 2, 3] :: Float
3.0

\end{verbatim}}
\end{quote}
Users should take care to pick a return type that is wide enough to contain
 the full length of the list. If the width is insufficient, the overflow
 behaviour will depend on the \haddocktt{(+)} implementation in the selected \haddockid{Num}
 instance. The following example overflows because the actual list length
 of 200 lies outside of the \haddocktt{Int8} range of \haddocktt{-128..127}.\par
\begin{quote}
{\haddockverb\begin{verbatim}
>>> genericLength [1..200] :: Int8
-56

\end{verbatim}}
\end{quote}}
\end{haddockdesc}
\begin{haddockdesc}
\item[\begin{tabular}{@{}l}
genericTake :: Integral i => i -> {\char 91}a{\char 93} -> {\char 91}a{\char 93}
\end{tabular}]
{\haddockbegindoc
The \haddockid{genericTake} function is an overloaded version of \haddockid{take}, which
 accepts any \haddockid{Integral} value as the number of elements to take.\par}
\end{haddockdesc}
\begin{haddockdesc}
\item[\begin{tabular}{@{}l}
genericDrop :: Integral i => i -> {\char 91}a{\char 93} -> {\char 91}a{\char 93}
\end{tabular}]
{\haddockbegindoc
The \haddockid{genericDrop} function is an overloaded version of \haddockid{drop}, which
 accepts any \haddockid{Integral} value as the number of elements to drop.\par}
\end{haddockdesc}
\begin{haddockdesc}
\item[\begin{tabular}{@{}l}
genericSplitAt :: Integral i => i -> {\char 91}a{\char 93} -> ({\char 91}a{\char 93}, {\char 91}a{\char 93})
\end{tabular}]
{\haddockbegindoc
The \haddockid{genericSplitAt} function is an overloaded version of \haddockid{splitAt}, which
 accepts any \haddockid{Integral} value as the position at which to split.\par}
\end{haddockdesc}
\begin{haddockdesc}
\item[\begin{tabular}{@{}l}
genericIndex :: Integral i => {\char 91}a{\char 93} -> i -> a
\end{tabular}]
{\haddockbegindoc
The \haddockid{genericIndex} function is an overloaded version of \haddockid{!!}, which
 accepts any \haddockid{Integral} value as the index.\par}
\end{haddockdesc}
\begin{haddockdesc}
\item[\begin{tabular}{@{}l}
genericReplicate :: Integral i => i -> a -> {\char 91}a{\char 93}
\end{tabular}]
{\haddockbegindoc
The \haddockid{genericReplicate} function is an overloaded version of \haddockid{replicate},
 which accepts any \haddockid{Integral} value as the number of repetitions to make.\par}
\end{haddockdesc}